% ---------------------------------------------------------------------------
% ШВИДКА КОМПІЛЯЦІЯ: під час написання тексту розкоментуйте рядок нижче —
% картинки замінюються рамками, PDF збирається в рази швидше.
% \PassOptionsToPackage{draft}{graphicx}
% ---------------------------------------------------------------------------
\documentclass[a4paper, 12pt, titlepage]{article}
\renewcommand{\baselinestretch}{1.5}

% --- Шрифти, мова, сторінка ---
\usepackage[T2A]{fontenc}                 % кирилиця (inputenc utf8 вже за замовчуванням)
\usepackage[english,ukrainian]{babel}
\usepackage[left=2.5cm,right=1cm,top=2cm,bottom=2cm]{geometry}
\usepackage{indentfirst}
\usepackage{fancyhdr}

% --- Математика ---
\usepackage{amsmath, amssymb, amsthm, mathtools, mathrsfs}
\DeclareMathOperator*{\argmax}{argmax}
\DeclareMathOperator*{\argmin}{argmin}
\newcommand{\RomanNumeralCaps}[1]{\MakeUppercase{\romannumeral #1}}

% --- Графіка, рисунки, таблиці ---
\usepackage{graphicx}
\usepackage{xcolor}
\usepackage{tikz}
\usetikzlibrary{arrows.meta,positioning,babel}
\usepackage{float}
\usepackage{placeins}
\usepackage{caption}
\usepackage{subcaption}
\usepackage{booktabs}
\usepackage{multirow}
\usepackage{array}

% --- Код (потрібно для jupyter.tex) ---
\usepackage{fvextra}
\usepackage{tcolorbox}
\usepackage{listings}
\tcbuselibrary{listingsutf8}

% --- Інше ---
\usepackage{dirtytalk}
\usepackage{comment}
\usepackage[shortlabels]{enumitem}
\usepackage[toc,page]{appendix}
\usepackage{etoolbox}
\patchcmd{\thebibliography}{\section*{\refname}}{}{}{}
\usepackage[hidelinks]{hyperref}          % hyperref — останнім

\addto\captionsukrainian{%
  \renewcommand{\contentsname}{\centering ЗМІСТ}%
}

\makeatletter
\renewcommand\appendix{\par
  \setcounter{section}{0}%
  \setcounter{subsection}{0}%
  \gdef\thesection{\@Asbuk\c@section}}
\makeatother

\input{jupyter}




\begin{document}


\input{title}
\pagestyle{myheadings}



\newpage
\thispagestyle{empty}
\tableofcontents



\newpage
\pagenumbering{arabic}
\setcounter{page}{3}
\section*{\centering ПЕРЕЛІК УМОВНИХ ПОЗНАЧЕНЬ}
\addcontentsline{toc}{section}{Перелік умовних позначень}
\noindent\textbf{БпЛА} -- безпілотний літальний апарат.

\noindent\textbf{GPS (Global Positioning System)} -- сукупність радіоелектронних засобів, що дає змогу визначати положення та швидкість руху об'єкта на поверхні Землі або в атмосфері.

\noindent\textbf{CNN (Convolution Neural Network)} -- це спеціалізований клас нейронних мереж, призначених для обробки сіткових даних, наприклад зображень. Вони особливо добре підходять для задач розпізнавання та обробки зображень.

\noindent\textbf{SLAM (Simultaneous Localization and Mapping)} -- метод одночасного навігації (локалізації) та побудови карти, який дозволяє автономному апарату будувати модель невідомого середовища і водночас визначати своє положення в ньому.

\noindent\textbf{SIFT (Scale-Invariant Feature Transform)} -- класичний алгоритм комп'ютерного зору для виявлення та опису локальних ознак зображення, що забезпечує високу інваріантність (стійкість) до змін масштабу, повороту та освітлення сцени.

\noindent\textbf{SURF (Speeded-Up Robust Features)} -- класичний детектора та дескриптор локальних ознак, що є прискореним аналогом методу SIFT за рахунок використання інтегральних зображень та апроксимації фільтрів Гаусса.

\noindent\textbf{ORB (Oriented FAST and Rotated BRIEF)} -- швидкий класичний алгоритм виділення ключових точок та формування бінарних дескрипторів, оптимізований для обчислень у реальному часі на пристроях із обмеженою потужністю.

\noindent\textbf{DISK (Learning Local Features with Policy Gradient)} -- сучасний нейромережевий метод екстракції локальних ознак, який навчається за допомогою алгоритмів підкріплення для досягнення стійкості до радикальних змін ракурсу та перспективи.



\newpage
\section*{\centering ВСТУП}
\addcontentsline{toc}{section}{Вступ}

Стрімкий розвиток безпілотних літальних апаратів (БпЛА) суттєво розширив сфери їхнього практичного застосування. Сучасні безпілотні системи активно використовуються у військовій сфері, моніторингу територій, картографуванні, рятувальних операціях, інспекції інфраструктури, логістиці та сільському господарстві. Для ефективного виконання таких задач критично важливим є забезпечення точної та стабільної навігації апарата у просторі.

Традиційно для визначення координат БпЛА використовуються глобальні навігаційні супутникові системи (GNSS/GPS). Проте в реальних умовах експлуатації сигнал GPS може бути нестабільним або повністю відсутнім через радіоелектронні перешкоди, навмисне приглушення сигналу, складний рельєф місцевості, щільну міську забудову або роботу в закритих просторах. У таких умовах особливої актуальності набувають методи візуальної локалізації, які базуються на аналізі зображень із бортових камер.

Однією з ключових задач комп’ютерного зору у системах візуальної навігації є зіставлення зображень. Суть цієї задачі полягає у пошуку відповідностей між поточним кадром, отриманим із камери БпЛА, та еталонною картою місцевості або ортофотопланом. Якість роботи алгоритмів зіставлення безпосередньо впливає на точність визначення координат апарата, стабільність траєкторії польоту та надійність автономної навігації.

Класичні методи комп’ютерного зору, такі як SIFT, SURF та ORB, тривалий час залишалися основними підходами до пошуку локальних ознак та встановлення відповідностей між зображеннями. Однак їхня ефективність суттєво знижується в умовах сильних змін масштабу, ракурсу, освітлення або за наявності слабкоструктурованих ділянок місцевості. Розвиток глибокого навчання та трансформерних архітектур дозволив створити нові нейромережеві методи, здатні забезпечувати значно вищу точність та стійкість зіставлення.

Особливий інтерес для задач візуальної локалізації БпЛА становлять сучасні методи LightGlue та LoFTR. Метод LightGlue використовує трансформерну архітектуру та адаптивні механізми скорочення обчислень, що дозволяє забезпечити високу швидкодію під час роботи в режимі реального часу. Метод LoFTR реалізує бездетекторний (\textit{detector-free}) підхід і виконує пошук відповідностей без явного виділення ключових точок, що дозволяє ефективно працювати навіть у слабкоструктурованих або однорідних областях зображення.

\textbf{\textit{Актуальність теми.}} Актуальність теми зумовлена необхідністю підвищення автономності та надійності навігаційних систем БпЛА в умовах нестабільної роботи або повної відсутності GPS. Сучасні методи зіставлення зображень на основі глибокого навчання відкривають нові можливості для побудови високоточних систем візуальної локалізації, здатних працювати у складних умовах навколишнього середовища.

Порівняльний аналіз сучасних трансформерних методів зіставлення дозволяє визначити їхні переваги, недоліки та доцільність використання у системах автономної навігації БпЛА реального часу.

\textbf{\textit{Мета дослідження.}} Метою магістерської роботи є розроблення системи візуальної локалізації безпілотних літальних апаратів за ортофотопланом на основі трансформерних методів зіставлення зображень LightGlue та LoFTR, експериментальне порівняння цих методів у складі системи та вдосконалення системи за результатами дослідження.

Для досягнення поставленої мети необхідно виконати такі \textbf{\textit{завдання}}:
\begin{enumerate}[label=\arabic*)]
    \item проаналізувати сучасні методи комп’ютерного зору для задач візуальної локалізації БпЛА;
    \item дослідити принципи роботи алгоритмів LightGlue та LoFTR;
    \item реалізувати програмну систему зіставлення кадрів БпЛА з ортофотопланом місцевості;
    \item провести експериментальне дослідження ефективності методів;
    \item виконати порівняльний аналіз точності та швидкодії алгоритмів;
    \item розробити та експериментально перевірити вдосконалення системи: вибір параметрів розбиття карти, пошук у вікні карти за прогнозом, уточнення положення, урахування очікуваного масштабу та згладжування траєкторії;
    \item перевірити систему на іншому типі місцевості з еталоном за даними супутникової навігації та дослідити використання курсу з бортового компаса;
    \item оцінити придатність досліджуваних методів для використання у системах автономної навігації БпЛА.
\end{enumerate}

\textbf{\textit{Об’єкт дослідження}} --- процес візуальної локалізації безпілотних літальних апаратів.

\textbf{\textit{Предмет дослідження}} --- сучасні методи зіставлення зображень на основі трансформерних архітектур для визначення положення БпЛА за даними бортової камери.

Робота складається зі вступу, трьох розділів, висновків та списку використаних джерел.

У \textit{першому розділі} розглянуто теоретичні основи візуальної локалізації БпЛА, принципи роботи сучасних методів комп’ютерного зору, механізми формування локальних ознак, трансформерні архітектури та алгоритми LightGlue і LoFTR.

У \textit{другому розділі} описано програмну реалізацію системи локалізації, процес підготовки ортофотоплану, алгоритм зіставлення кадрів БпЛА з картою, особливості реалізації досліджуваних методів, а також запропоновані вдосконалення системи.

У \textit{третьому розділі} наведено результати експериментального дослідження, виконано порівняльний аналіз точності та швидкодії методів LightGlue і LoFTR, оцінено ефект запропонованих вдосконалень, перевірено систему на другому польоті над сільськогосподарською місцевістю, зокрема з використанням курсу з компаса, а також сформульовано висновки щодо ефективності методів у задачах візуальної локалізації БпЛА.


\newpage


\section{Теоретичні основи візуальної локалізації БпЛА}
\subsection{Задача абсолютної візуальної навігації літальних апаратів}

У реальних умовах застосування БпЛА часто стикаються з нестабільністю або повною відсутністю сигналів GNSS/GPS через радіоелектронне придушення (РЕБ) чи фізичні перешкоди \cite{gps}. У таких випадках основним джерелом навігаційної інформації стає бортова камера \cite{visual_nav}.

На відміну від візуальної одометрії (VO) та методів SLAM, які обчислюють відносне зміщення і неминуче накопичують похибку (дрейф) з часом, \textbf{абсолютна візуальна локалізація} зіставляє поточний кадр із заздалегідь підготовленим ортофотопланом місцевості. Це забезпечує точну прив'язку координат без накопичувальної похибки.

Математично задача зводиться до пошуку фрагмента карти та матриці просторового перетворення (гомографії) $H$, які найкраще геометрично узгоджують кадр із картою. Проте на практиці цей процес стикається з низкою критичних викликів:

\begin{itemize}
    \item \textbf{Масштаб та просторова дискретизація:} Кадр охоплює лише малу частину глобального ортофотоплану з іншою просторовою роздільною здатністю, що вимагає нормалізації масштабу зображень та розбиття карти на дрібніші пошукові фрагменти (тайли).
    \item \textbf{Міжсезонні та фотометричні варіації:} Зйомка часто відбувається в умовах (наприклад, узимку), що кардинально відрізняються від умов створення карти (улітку). Це призводить до зміни тіней і текстур, що руйнує традиційні піксельні метрики подібності та вимагає вилучення високорівневих структурних ознак сцени.
    \item \textbf{Паралакс рельєфу:} Модель проєктивної гомографії коректно описує лише пласкі поверхні. Висотні об'єкти (багатоповерхові будівлі) створюють проєктивні зміщення, через які алгоритм може хибно прив'язуватися до площини даху замість землі.
    \item \textbf{Неінваріантність до курсу (In-plane rotation):} Сучасні нейромережеві методи мають слабку стійкість до сильного обертання зображення (деградують при кутах понад $30^\circ$--$45^\circ$). Під час маневрування БпЛА це вимагає явної компенсації курсу перед етапом зіставлення.
\end{itemize}

Саме ці фактори роблять класичні алгоритми комп'ютерного зору недостатньо надійними, зумовлюючи перехід до сучасних трансформерних архітектур, здатних встановлювати відповідності у складних умовах змінного середовища.

\subsection{Еволюція виділення та опису візуальних ознак}

Традиційно базою для систем комп'ютерного зору слугували класичні детектори та дескриптори локальних ознак, такі як SIFT \cite{sift} та ORB \cite{orb}. Їхня математична основа (аналіз градієнтів та інтенсивностей) забезпечує високу стійкість до афінних перетворень \cite{classic_methods}. Однак під час вирішення задач навігації БпЛА ці методи виявилися вразливими до сильних фотометричних та міжсезонних змін сцени. Оскільки вони спираються на локальну піксельну статистику, зміна снігового покриву, кольору крон дерев або довжини тіней призводить до неможливості встановлення надійних відповідностей.

Подолання цих обмежень стало можливим завдяки переходу до нейромережевих методів на основі згорткових нейронних мереж (CNN). Сучасні архітектури, такі як SuperPoint \cite{superpoint}, виконують роль комплексних кодувальників геометрії (\textit{geometry encoders}). Замість ручного проєктування математичних фільтрів, вони за допомогою глибокого навчання самостійно формують високорівневі семантичні абстракції об'єктів, що забезпечує значно вищу стійкість виділення ключових точок (\textit{keypoints}) та їхніх векторів-дескрипторів у складних умовах навколишнього середовища.

Для розвитку бездетекторних (\textit{detector-free}) систем зіставлення еволюція пішла шляхом повної відмови від виділення дискретних точок на користь щільних карт ознак (\textit{dense feature maps}). Згорткові кодувальники на базі залишкових мереж (ResNet) \cite{resnet} у поєднанні з архітектурою пірамід ознак (FPN) \cite{cnn_fpn} формують ієрархічне просторове представлення зображення. Цей підхід забезпечує генерацію двох типів карт:
\begin{itemize}
    \item \textbf{Грубі ознаки (Coarse level):} карти низької просторової роздільної здатності (масштаб $1/8$ від вихідного зображення). Вони охоплюють широке рецептивне поле та кодують глобальний семантичний контекст сцени, що дозволяє знаходити відповідні ділянки навіть на слабкотекстурованих поверхнях (лісові масиви, поля).
    \item \textbf{Точні ознаки (Fine level):} карти високої просторової роздільної здатності (масштаб $1/2$), що зберігають чіткі локальні межі та текстурні переходи. Вони використовуються для фінального субпіксельного уточнення просторових координат.
\end{itemize}

Ця стратегія переходу від загального до часткового (\textit{coarse-to-fine}) стала математичним фундаментом для сучасних нейромережевих конвеєрів зіставлення, дозволяючи оптимізувати обчислювальну складність наступних трансформерних шарів.

\subsection{Трансформерні архітектури встановлення відповідностей}

Наступним етапом після вилучення візуальних ознак є пошук просторових відповідностей між ними. Сучасним стандартом для цього завдання стали архітектури на основі трансформерів \cite{attention_all_need}. Їхньою основою є механізми уваги, які дозволяють встановлювати залежності між елементами незалежно від просторової відстані між ними на зображенні, формуючи глибокий глобальний контекст сцени.

\subsubsection{Механізми внутрішньої та перехресної уваги}
У задачах візуальної локалізації конвеєр уваги традиційно розділяють на два етапи, що чергуються:
\begin{itemize}
    \item \textbf{Внутрішня увага (\textit{Self-Attention}):} аналізує геометричні та структурні зв'язки між ознаками всередині одного зображення. Завдяки цьому кожен локальний дескриптор оновлюється з урахуванням глобального контексту (наприклад, точка «розуміє», що вона належить до кута будинку, який стоїть поруч із перехрестям).
    \item \textbf{Перехресна увага (\textit{Cross-Attention}):} безпосередньо порівнює ознаки поточного кадру БпЛА з ознаками фрагмента карти, формуючи інформаційні зв'язки між потенційно відповідними візуальними структурами обох зображень.
\end{itemize}

Залежно від того, які саме ознаки подаються на вхід трансформерним блокам, виділяють детектор-залежні (розріджені) та бездетекторні (щільні) методи зіставлення.

\subsubsection{Метод розрідженого зіставлення LightGlue}
Метод LightGlue \cite{lightglue} розроблено для швидкого та високоточного зіставлення розріджених наборів дискретних ключових точок. Вхідними даними є координати $p^A, p^B$ та локальні дескриптори $d^A, d^B$, сформовані зовнішнім екстрактором (наприклад, SuperPoint \cite{superpoint}). На нульовому шарі мережі стан ознак ініціалізується виключно проєкцією нормалізованих векторів дескрипторів:
\begin{equation}
x_i^{(0)} = W_{\text{in}} d_i + b_{\text{in}}.
\end{equation}

Геометричні координати точок $p_i = (u_i, v_i)$ не змішуються з дескриптором, а враховуються у блоках внутрішньої уваги за допомогою обертального позиційного кодування (\textit{Rotary Position Embedding, RoPE}):
\begin{equation}
\text{Attn}(Q, K, V) = \text{softmax}\left( \frac{\mathcal{R}(Q, p) \cdot \mathcal{R}(K, p)^T}{\sqrt{d_k}} \right) V,
\end{equation}
де $\mathcal{R}(\cdot, p)$ --- оператор двовимірного обертання підпросторів векторів запитів і ключів. Такий підхід забезпечує залежність скалярного добутку уваги виключно від відносної відстані між точками, зберігаючи трансляційну інваріантність.

\textbf{Динамічне відсікання ознак та адаптивна зупинка.}\\
Ефективність LightGlue досягається за рахунок динамічної адаптації обчислювальної складності. На кожному шарі $l$ класифікатор прогнозує здатність точки мати надійну відповідність (\textit{matchability}): $c_i^{(l)} = \sigma( \text{MLP}_{\text{conf}}( x_i^{(l)} ) )$. Це дозволяє реалізувати дві оптимізації:
\begin{enumerate}
    \item \textbf{Відсікання (\textit{Point Pruning}):} якщо $c_i^{(l)} < \tau_{\text{prune}}$, точка маркується як неінформативна (наприклад, потрапила за межі спільного огляду або на монотонну поверхню) і повністю вилучається з подальшої обробки, зменшуючи розмірність матриць уваги.
    \item \textbf{Рання зупинка (\textit{Early Stopping}):} якщо достатня частка точок досягла високого рівня впевненості ($c_i^{(l)} > \tau_{\text{stop}}$), конвеєр обчислень достроково припиняє виконання наступних шарів трансформера.
\end{enumerate}

\textbf{Фінальне призначення відповідностей.}\\
На виході з останнього шару формується матриця подібності $S$. Пара точок $(i, j)$ визнається коректною відповідністю, якщо виконується строгий критерій взаємних найближчих сусідів (\textit{Mutual Nearest Neighbors, MNN}) та ймовірність збігу перевищує заданий поріг фільтрації:
\begin{equation}
P_{i,j} > \tau_{\text{match}}, \quad \text{де } j = \argmax_{k} S_{i,k} \quad \text{та} \quad i = \argmax_{m} S_{m,j}.
\end{equation}

\subsubsection{Бездетекторне щільне зіставлення LoFTR}
На відміну від LightGlue, метод LoFTR \cite{loftr} повністю відмовляється від попереднього виділення ключових точок. Спираючись на ієрархічні карти ознак згорткового кодувальника (див. підрозділ 1.2), LoFTR виконує зіставлення безпосередньо між семантичними елементами сцени.

\textbf{Грубе зіставлення з механізмом Dual-Softmax.}\\
Отримані грубі карти ознак $\tilde{F}_A$ та $\tilde{F}_B$ (масштаб $1/8$) розгортаються в одновимірні послідовності, доповнюються абсолютним позиційним кодуванням і обробляються блоками внутрішньої та перехресної уваги. Далі обчислюється матриця подібності $S$ на основі скалярного добутку трансформованих ознак:
\begin{equation}
S(i, j) = \frac{1}{\tau} \langle \tilde{F}_A(i), \tilde{F}_B(j) \rangle.
\end{equation}
Для перетворення сирих значень матриці $S$ у матрицю ймовірностей збігу $P_c$ застосовується операція двостороннього софтмаксу (\textit{Dual-Softmax}):
\begin{equation}
P_{c}(i, j) = \text{softmax}(S(i, \cdot))_j \cdot \text{softmax}(S(\cdot, j))_i.
\end{equation}
Цей підхід ефективно пригнічує хибні відповідності на монотонних або повторюваних структурах (однотипні дахи будівель, крони дерев), оскільки наявність кількох схожих кандидатів геометрично збільшує знаменник софтмаксу та знижує фінальну ймовірність $P_c$ для кожного з них.

\textbf{Точне субпіксельне уточнення.}\\
Для відібраних грубих пар точок алгоритм переходить до карт ознак високої роздільної здатності $\hat{F}_A, \hat{F}_B$ (масштаб $1/2$). Навколо координат грубого збігу вирізаються локальні вікна розміром $w \times w$ пікселів, де знову застосовуються невеликі трансформерні шари.

Замість жорсткого вибору дискретного пікселя з найбільшою ймовірністю, фінальна координата $\hat{j}$ розраховується як математичне сподівання просторового розподілу ймовірностей у межах локального вікна:
\begin{equation}
\hat{j} = \sum_{n \in w \times w} P(n) \cdot \mathbf{x}_n,
\end{equation}
де $\mathbf{x}_n$ --- координати пікселя $n$ у вікні, $P(n)$ --- нормована ймовірність схожості. Таке субпіксельне уточнення є критично важливим для систем візуальної навігації БпЛА, адже помилка в один піксель на кадрах, знятих із великої висоти, трансформується у суттєве зміщення реальних координат апарата.

\subsection{Математичні моделі просторового перетворення та робастного оцінювання}

Встановлення множини точкових відповідностей між кадром та фрагментом карти є лише першим етапом локалізації. Оскільки отриманий набір пар неминуче містить хибні збіги (\textit{outliers}), викликані візуальними артефактами або повторюваними текстурами, для визначення істинних координат камери необхідно побудувати геометрично узгоджену математичну модель перетворення.

\subsubsection{Проєктивна гомографія та декомпозиція якобіана}
За умови достатньої висоти польоту БпЛА та надирного напрямку зйомки перепад висот рельєфу є малим порівняно з відстанню до камери. Це дозволяє наближено вважати зняту сцену площиною. Зв'язок між координатами точок на двох площинах (кадрі та карті) у комп'ютерному зорі повністю описується матрицею проєктивної гомографії $H \in \mathbb{R}^{3 \times 3}$ \cite{hartley}.

Якщо точка на бортовому кадрі має однорідні координати $\mathbf{p} = [x, y, 1]^T$, а відповідна їй точка на еталонному ортофотоплані -- $\mathbf{p}' = [x', y', 1]^T$, їхнє проєктивне перетворення виражається системою:
\begin{equation}
s \begin{bmatrix} x' \\ y' \\ 1 \end{bmatrix} = H \begin{bmatrix} x \\ y \\ 1 \end{bmatrix} = \begin{bmatrix} h_{11} & h_{12} & h_{13} \\ h_{21} & h_{22} & h_{23} \\ h_{31} & h_{32} & h_{33} \end{bmatrix} \begin{bmatrix} x \\ y \\ 1 \end{bmatrix},
\label{eq:homography_matrix}
\end{equation}
де $s$ -- довільний ненульовий масштабувальний коефіцієнт. Оскільки матриця $H$ визначена з точністю до масштабу, вона має 8 ступенів вільності і потребує мінімум 4 неколінеарних пар точок для розв'язання системи лінійних рівнянь. Перехід від однорідних координат до фізичних (декартових) на площині карти здійснюється через операцію дегомогенізації:
\begin{equation}
\tilde{\mathbf{p}}' = H \mathbf{p} = \begin{bmatrix} \tilde{x}' \\ \tilde{y}' \\ \tilde{w}' \end{bmatrix}, \qquad x' = \frac{\tilde{x}'}{\tilde{w}'}, \quad y' = \frac{\tilde{y}'}{\tilde{w}'}.
\end{equation}

Теоретично алгоритми робастного оцінювання можуть знайти формально правильну матрицю $H$, яка мінімізує похибку на підмножині точок, але описує фізично неможливе перетворення (наприклад, сильне зминання або дзеркальне вивертання площини). Для перевірки правдоподібності знайденої моделі застосовується \textbf{локальне афінне наближення гомографії}. Обчислюючи матрицю Якобі $J$ проєктивного відображення у фіксованій точці (зазвичай у центрі кадру $\mathbf{c}$), можна розкласти перетворення на зрозумілі фізичні компоненти:
\begin{align}
\sigma &= \sqrt{|\det J|} && \text{--- масштаб (відношення площ)}, \\
\theta &= \operatorname{atan2}(J_{21} - J_{12},\, J_{11} + J_{22}) && \text{--- кут повороту (курс БпЛА відносно карти)}, \\
\kappa &= \frac{\sigma_{\max}(J)}{\sigma_{\min}(J)} && \text{--- анізотропія (деформація пропорцій)},
\end{align}
де $\sigma_{\max}$ та $\sigma_{\min}$ -- сингулярні числа якобіана. Це дозволяє впроваджувати жорсткі фізичні обмеження на етапі валідації гомографії, такі як вимога збереження орієнтації ($\det J > 0$) та обмеження на максимальну анізотропію ($\kappa \le 1{,}5$ для надирної камери). Крім того, наявність висотних об'єктів (наприклад, будівель) спричиняє паралакс рельєфу. У таких випадках матриця $H$ може бути помилково оцінена за площиною даху, а не землі, що вимагає ретельного контролю розрахованого масштабу $\sigma$.

\subsubsection{Алгоритми робастного консенсусу: від RANSAC до MAGSAC++}
Оскільки множина вхідних відповідностей засмічена хибними парами, пряме розв'язання рівняння (\ref{eq:homography_matrix}) методом найменших квадратів неможливе. Класичним рішенням є ітераційний алгоритм RANSAC \cite{ransac}, який випадковим чином обирає мінімальні вибірки (по 4 пари), будує для кожної з них гіпотезу матриці $H$ та підраховує кількість точок підтримки (\textit{inliers}), чия похибка репроєкції не перевищує фіксований жорсткий поріг $\tau$.

Попри свою популярність, класичний RANSAC має два суттєві недоліки для систем реального часу. По-перше, вибір фіксованого порогу $\tau$ є компромісом: замалий поріг відкидає корисні зашумлені точки, а завеликий -- пропускає хибні гіпотези. По-друге, за відсутності істинної моделі (наприклад, коли кадр порівнюється з фрагментом карти, який не має з ним геометричного перетину) класичний RANSAC марно виконує всі задані ітерації (зазвичай тисячі), витрачаючи значний обчислювальний час.

Алгоритм \textbf{MAGSAC++} (\textit{Marginalizing Sample Consensus}) \cite{magsac} усуває ці недоліки, відмовляючись від фіксованого порогу відсікання. Він розглядає рівень шуму як неперервну випадкову величину $\sigma$ із заданим апостеріорним розподілом і маргіналізує (інтегрує) похибку за всіма можливими значеннями дисперсії:
\begin{equation}
L(e_i) = \int_0^{\sigma_{\max}} P(e_i \mid \sigma) P(\sigma) \, d\sigma,
\end{equation}
де $e_i$ -- похибка репроєкції $i$-ї точки. Замість бінарної оцінки кожна точка отримує неперервну вагу підтримки гіпотези, яка плавно спадає зі зростанням невязки.

Окрім підвищення точності, застосування MAGSAC++ суттєво пришвидшує роботу системи. Завдяки вбудованим механізмам пре-емптивного відсікання (\textit{Sequential Probability Ratio Test, SPRT}), алгоритм розпізнає безперспективні гіпотези на ранніх етапах оцінювання і миттєво перериває ітерації для пар зображень, що не мають перекриття. Це є критично важливим для локалізації в просторих вікнах пошуку, зменшуючи обчислювальні витрати в кілька разів без втрати надійності геометричної моделі.

 та кутів орієнтації (крен, тангаж, рискання) апарата у просторі.

Типовий конвеєр обробки даних системи візуальної локалізації містить п'ять послідовних етапів:

\begin{center}
\begin{minipage}{0.6\textwidth}
\ttfamily
\centering
Реєстрація зображення \\
$\downarrow$ \\
Виділення ознак \\
$\downarrow$ \\
Зіставлення ознак \\
$\downarrow$ \\
Геометрична верифікація \\
$\downarrow$ \\
Оцінювання пози БпЛА
\end{minipage}
\end{center}
\vspace{0.3cm}

% \begin{itemize}
%     \item Реєстрація зображення: фіксація оптичного сигналу матрицею бортової камери, попереднє цифрове фільтрування шумів та нормалізація кадру.
%     \item \textbf{Виділення ознак (\textit{Feature extraction}):} виявлення стабільних локальних орієнтирів (ключових точок) та обчислення їхніх математичних дескрипторів (або генерація щільних карт ознак у бездетекторних методах типу LoFTR).
%     \item \textbf{Зіставлення ознак (\textit{Feature matching}):} пошук взаємно однозначних пар відповідностей між векторами дескрипторів бортового кадру та еталонного фрагмента ортофотоплану.
%     \item \textbf{Геометрична верифікація (\textit{Geometric verification}):} застосування робастних статистичних методів (наприклад, RANSAC) для фільтрації помилкових збігів та оцінювання математичної матриці плоскої гомографії.
%     \item \textbf{Оцінювання пози (\textit{Pose estimation}):} розв'язання системи проективних рівнянь (зокрема, задачі Perspective-n-Point, PnP) для точного відновлення тривимірних координат і векторів кутової орієнтації камери відносно площини карти.
% \end{itemize}

\textbf{Основні проблеми візуальної локалізації}\\
Ефективність функціонування розробленого конвеєра у реальних місіях БпЛА суттєво ускладнюється наявністю п'яти основних фізичних та геометричних проблем:
\begin{itemize}
    \item Зміни умов освітлення: коливання інтенсивності світла протягом доби, пряме сонячне світло, утворення довгих контрастних тіней від будівель та дерев, які спотворюють математичні значення дескрипторів.
    \item Радикальні зміни ракурсу: сильні перспективні спотворення та зміни форми об'єктів, викликані нахилами БпЛА під час маневрування, а також розбіжності у кутах зйомки між бортовою камерою та ортофотопланом (який зазвичай знятий суворо ортогонально вниз).
    \item Ефект змазування зображення: розмиття текстури кадру внаслідок високошвидкісного руху апарата або кутових вібрацій його корпусу, що унеможливлює роботу класичних детекторів точок.
    \item Сезонні та кліматичні варіації: кардинальна зміна візуального вигляду однієї й тієї самої місцевості в різні пори року (наявність снігового покриву, цвітіння або опадіння листя дерев), що руйнує стабільність класичних ознак.
    \item Слабкотекстуровані області: наявність на маршруті польоту протяжних візуально однорідних поверхонь (водні плеса, піщані зони, засніжені рівнини, монотонні дахи промислових комплексів), де класичні алгоритми не здатні виділити унікальні ключові точки.
\end{itemize}
Саме для подолання зазначених проблем у даній роботі досліджуються сучасні трансформерні архітектури LightGlue та LoFTR, які використовують контекстні механізми уваги для забезпечення надійності зіставлення.

\subsection{Основи рекурсивного оцінювання стану та фільтрації траєкторій}

Методи геометричного зіставлення генерують незалежні оцінки координат БпЛА для кожного окремого кадру. Оскільки такі вимірювання неминуче містять випадкові флуктуації (шум спостережень), а іноді можуть бути тимчасово недоступними (через розмиття кадру або відсутність текстури), для побудови безперервної та гладкої траєкторії польоту застосовуються методи рекурсивного баєсівського оцінювання.

Основою таких методів є модель простору станів (\textit{state-space model}), яка описує кінематику літального апарата. Для задач планарної навігації стан системи в дискретний момент часу $k$ найчастіше описується вектором $\mathbf{x}_k = [X_k, Y_k, v_{X,k}, v_{Y,k}]^T$, що містить координати та складові швидкості. Динаміка системи (наприклад, модель руху зі сталою швидкістю) та процес вимірювання описуються системою лінійних стохастичних рівнянь:
\begin{align}
\mathbf{x}_k &= F \mathbf{x}_{k-1} + \mathbf{w}_k, \\
\mathbf{z}_k &= H \mathbf{x}_k + \mathbf{v}_k,
\end{align}
де $F$ --- матриця переходу стану, $H$ --- матриця спостережень, яка проєктує повний вектор стану на простір вимірюваних координат (у нашому випадку --- на візуально оцінені координати $[X, Y]^T$). Системний шум $\mathbf{w}_k \sim \mathcal{N}(0, Q)$ моделює невраховані прискорення (маневри БпЛА, пориви вітру), а шум вимірювань $\mathbf{v}_k \sim \mathcal{N}(0, R)$ відображає похибку методу візуальної локалізації.

\subsubsection{Фільтр Калмана (Kalman Filter)}
Оптимальним розв'язком задачі рекурсивного оцінювання стану для лінійних систем із гауссовими шумами є дискретний фільтр Калмана \cite{kalman}. Робота фільтра складається з двох кроків, що чергуються в реальному часі:
\begin{enumerate}
    \item \textbf{Етап прогнозу (\textit{Predict}):} система екстраполює вектор стану та матрицю його коваріації $P$ на наступний крок за допомогою кінематичної моделі:
    \begin{align}
    \hat{\mathbf{x}}_{k|k-1} &= F \hat{\mathbf{x}}_{k-1|k-1}, \\
    P_{k|k-1} &= F P_{k-1|k-1} F^T + Q.
    \end{align}
    \item \textbf{Етап корекції (\textit{Update}):} щойно надходить нове візуальне вимірювання $\mathbf{z}_k$, фільтр обчислює невязку (\textit{innovation}) та матрицю коефіцієнтів підсилення Калмана $K_k$. Оцінка стану коригується пропорційно до цього підсилення:
    \begin{equation}
    \hat{\mathbf{x}}_{k|k} = \hat{\mathbf{x}}_{k|k-1} + K_k (\mathbf{z}_k - H \hat{\mathbf{x}}_{k|k-1}).
    \end{equation}
\end{enumerate}
Фільтр Калмана використовує лише поточні та минулі вимірювання, що робить його придатним для роботи в режимі реального часу (\textit{online tracking}). Екстрапольовані оцінки $\hat{\mathbf{x}}_{k|k-1}$ дозволяють ефективно передбачати положення БпЛА для звуження вікна пошуку тайлів на наступному кадрі.

\subsubsection{Апостеріорне згладжування (RTS Smoother)}
Попри ефективність прямої фільтрації, оцінка $\hat{\mathbf{x}}_{k|k}$ не містить інформації про майбутні маневри апарата, тому відновлена траєкторія може мати певне запізнення в реакції на різкі зміни курсу. Для аналізу траєкторії після завершення польоту (\textit{offline analysis}) застосовують алгоритми згладжування з фіксованим інтервалом.

Стандартом для лінійних систем є згладжувач Рауха–Тунга–Штрібеля (RTS Smoother) \cite{rts}. Після того як прямий прохід фільтра Калмана зберігає всі апріорні та апостеріорні оцінки, згладжувач виконує зворотний прохід (\textit{backward pass}) від останнього кадру $N$ до першого:
\begin{align}
C_k &= P_{k|k} F^T P_{k+1|k}^{-1}, \\
\hat{\mathbf{x}}_{k|N} &= \hat{\mathbf{x}}_{k|k} + C_k (\hat{\mathbf{x}}_{k+1|N} - \hat{\mathbf{x}}_{k+1|k}),
\end{align}
де $\hat{\mathbf{x}}_{k|N}$ --- згладжена оцінка стану в момент $k$ з урахуванням \textit{усіх} вимірювань за час польоту. Застосування згладжувача RTS дозволяє суттєво зменшити коваріацію похибки та сформувати максимально плавну і точну траєкторію руху БпЛА, що є особливо корисним на ділянках із пропущеними або відхиленими візуальними вимірюваннями.




% =============================================================================
% Розділ 2. Програмне забезпечення та програмна реалізація системи
% Підключення: \input{sections/section2_implementation}
% Потрібно в преамбулі: \usetikzlibrary{arrows.meta,positioning,babel}
% =============================================================================
\providecommand{\px}{пкс}               % пікселі карти
\providecommand{\code}[1]{\texttt{#1}}  % імена файлів, модулів і функцій
\newpage
\section{Програмне забезпечення та програмна реалізація системи}
\label{sec:implementation}

\subsection{Постановка задачі та загальна архітектура системи}

Розроблена система розв'язує задачу абсолютної візуальної локалізації БпЛА: за кадром із бортової камери та ортофотопланом місцевості визначити, над якою точкою карти перебуває апарат, а за послідовністю кадрів --- відновити траєкторію польоту. Інформація про те, якому фрагменту карти відповідає кадр, системі не надається: відповідну ділянку ортофотоплану система знаходить автоматично.

Вхідними даними є:
\begin{itemize}
    \item ортофотоплан місцевості розміром $W \times H$ пікселів;
    \item відеозапис польоту БпЛА або послідовність кадрів із часовими мітками.
\end{itemize}

Вихідними даними є:
\begin{itemize}
    \item для кожного кадру --- координати $(X, Y)$ проєкції центра кадру на ортофотоплан і контур огляду камери (чотирикутник на карті);
    \item для послідовності кадрів --- траєкторія польоту з позначкою статусу кожного кадру та показниками якості зіставлення.
\end{itemize}

Робота системи ґрунтується на таких припущеннях:
\begin{itemize}
    \item камера спрямована приблизно вертикально вниз (надирна зйомка);
    \item сцену в межах кадру можна наближено вважати плоскою, тому перетворення <<кадр $\to$ карта>> описується гомографією \cite{hartley};
    \item відношення просторової роздільності кадру та карти відоме з точністю до множника $\approx 2$ (див. п.~\ref{sec:scale});
    \item усі координати обчислюються в пікселях карти (\px{}); географічна прив'язка ортофотоплану системі не потрібна. Ортофотоплан першого польоту прив'язки не має, тому похибки для нього наведено в пікселях; для другого польоту використано прив'язаний знімок, і похибки переведено в метри (п.~\ref{sec:datasets}).
\end{itemize}

Загальну схему обробки наведено на рисунку~\ref{fig:pipeline}. Система складається з офлайн-етапу підготовки карти, який виконується один раз, та онлайн-етапу обробки кадрів. Ключовим елементом онлайн-етапу є єдиний конвеєр зіставлення пари <<кадр --- фрагмент карти>>, який однаково застосовується в усіх режимах роботи: під час перевірки відомої пари, автоматичного пошуку фрагмента та трекінгу послідовності.

\begin{figure}[H]
\centering
\begin{tikzpicture}[
    node distance=5mm and 6mm,
    box/.style={draw, rounded corners=2pt, align=center, font=\small,
                minimum width=42mm, minimum height=9mm, inner sep=3pt},
    off/.style={box, fill=gray!12},
    on/.style={box, fill=blue!7},
    dec/.style={box, fill=orange!12},
    arr/.style={-{Stealth[length=2.2mm]}, thick}]
  % офлайн-етап
  \node[off] (map) {Ортофотоплан\\$6200\times4320$};
  \node[off, below=of map] (tiles) {Розбиття на фрагменти\\$512\times512$, перекриття 50\,\%};
  \node[off, below=of tiles] (cache) {Кеш ознак фрагментів\\(SuperPoint)};
  % онлайн-етап
  \node[on, right=16mm of map] (frame) {Кадр БпЛА\\$1920\times1080$};
  \node[on, below=of frame] (prep) {Нормалізація масштабу\\і компенсація курсу\\(прогноз або компас)};
  \node[dec, below=of prep] (pool) {Ділянки карти: уся карта,\\вікно за прогнозом / $3\times3$};
  \node[on, below=of pool] (match) {Зіставлення\\LightGlue / LoFTR};
  \node[on, below=of match] (geom) {RANSAC + гомографія,\\перевірка правдоподібності};
  \node[dec, below=of geom] (rank) {Фільтр і ранжування\\фрагментів};
  \node[on, below=of rank] (coord) {Координати $(X,Y)$\\і контур огляду};
  \node[dec, below=of coord] (track) {Кінематична перевірка,\\траєкторія};
  \node[off, below=of track] (smooth) {Згладжування траєкторії\\(фільтр Калмана)};
  \draw[arr] (track) -- (smooth);
  \draw[arr] (map) -- (tiles);
  \draw[arr] (tiles) -- (cache);
  \draw[arr] (frame) -- (prep);
  \draw[arr] (prep) -- (pool);
  \draw[arr] (pool) -- (match);
  \draw[arr] (match) -- (geom);
  \draw[arr] (geom) -- (rank);
  \draw[arr] (rank) -- (coord);
  \draw[arr] (coord) -- (track);
  \draw[arr] (cache.south) |- (match.west);
  \draw[arr, dashed] (track.east) -- ++(6mm,0) |- node[pos=0.25, right, font=\footnotesize, align=left]{прогноз\\положення\\і курсу} (pool.east);
\end{tikzpicture}
\caption{Загальна схема роботи системи візуальної локалізації (сірим позначено офлайн-етапи: підготовку карти та необов'язкове згладжування готової траєкторії)}
\label{fig:pipeline}
\end{figure}

\subsection{Використані програмні засоби}

Систему реалізовано мовою Python~3.11. Використано такі бібліотеки:
\begin{itemize}[label=\textbf{•}]
    \item \textbf{PyTorch} 2.14 з підтримкою CUDA~12.6 \cite{pytorch} --- виконання нейромережевих моделей на GPU;
    \item \textbf{LightGlue} (офіційна реалізація авторів) --- моделі SuperPoint \cite{superpoint} і LightGlue \cite{lightglue};
    \item \textbf{Kornia} 0.8 \cite{kornia} --- реалізація LoFTR \cite{loftr} з попередньо навченими вагами \textit{outdoor};
    \item \textbf{OpenCV} 5.0 \cite{opencv} --- читання та масштабування зображень, оцінювання гомографії методами RANSAC \cite{ransac} і MAGSAC++ \cite{magsac}, геометричні перетворення;
    \item \textbf{NumPy} і \textbf{pandas} --- числові обчислення, робота з таблицями метаданих і результатів;
    \item \textbf{Matplotlib} --- побудова графіків;
    \item \textbf{pytest} --- автоматичне тестування модулів.
\end{itemize}

Нейромережі в роботі не навчалися: використано опубліковані попередньо навчені ваги моделей.

\subsection{Тестове середовище}

Розроблення та експерименти виконувалися на ноутбуці з такими характеристиками:
\begin{itemize}[label=\textbf{--}]
    \item \textbf{процесор:} AMD Ryzen 5 4600H;
    \item \textbf{оперативна пам'ять:} 8 ГБ;
    \item \textbf{відеокарта:} NVIDIA GeForce GTX 1650 (4 ГБ відеопам'яті);
    \item \textbf{CUDA:} 12.6 (збірка PyTorch \texttt{cu126});
    \item \textbf{операційна система:} Windows 11.
\end{itemize}

\subsection{Вхідні дані}

Для моделювання використано ортофотоплан району розміром $6200 \times 4320$ пікселів (формат RGBA з повністю непрозорим альфа-каналом) і відеозапис польоту БпЛА DJI (файл \code{DJI\_0331.MP4}): роздільність $1920 \times 1080$, частота 29,97~кадр/с, тривалість $\approx 94$~с, кодек HEVC, камера спрямована вниз.

Відео та ортофотоплан знято в різні пори року: ортофотоплан --- влітку, відео --- взимку. Через це рослинність, тіні та покриття ґрунту на кадрах і на карті суттєво відрізняються, що ускладнює зіставлення.

Для перевірки системи на іншому типі місцевості використано другий набір даних --- політ над полями з геоприв'язаним супутниковим знімком і журналом бортової навігації (пп.~\ref{sec:datasets}, \ref{sec:exp-flight2}). Далі в цьому розділі параметри наведено для першого польоту.

\subsection{Підготовка ортофотоплану}

\textbf{Розбиття на фрагменти з перекриттям.}\\
Нейромережеві методи зіставлення розраховані на зображення розміром порядку сотень пікселів, тому ортофотоплан розбивається на квадратні фрагменти (тайли) розміром $S = 512$ пікселів із часткою перекриття $o = 0{,}5$. Крок сітки обчислюється за формулою
\begin{equation}
s = \lfloor S \cdot (1 - o) \rfloor = 256 .
\end{equation}
Координати лівих верхніх кутів фрагментів уздовж осі $x$ утворюють послідовність $0, s, 2s, \dots$; якщо регулярна сітка не доходить до правого краю карти, додається ще один крайовий фрагмент із початком у $W - S$ (аналогічно для осі $y$). Завдяки цьому всі фрагменти мають однаковий розмір і кожен піксель карти покрито. Кількість фрагментів:
\begin{equation}
N = \left( \left\lceil \frac{W - S}{s} \right\rceil + 1 \right) \cdot \left( \left\lceil \frac{H - S}{s} \right\rceil + 1 \right) = 24 \cdot 16 = 384 .
\end{equation}
Для кожного фрагмента в таблицю метаданих \code{tiles.csv} записуються ідентифікатор і координати $(X_0, Y_0)$ лівого верхнього кута в пікселях повної карти. Ці координати є основою для переходу від локальних координат фрагмента до глобальних координат карти.

Перекриття 50\,\% гарантує, що будь-яка ділянка карти, менша за половину фрагмента, повністю потрапляє принаймні в один фрагмент. Водночас, як показано в розділі~\ref{sec:experiments}, контур огляду камери на карті має розмір $\approx 518 \times 292$~\px{} і є ширшим за фрагмент, тому кадр, як правило, лише частково перекривається з кожним окремим фрагментом.

\textbf{Альтернативне розбиття.}\\
Розмір фрагмента та частка перекриття є параметрами системи. Для дослідження їхнього впливу (п.~\ref{sec:exp-improve}) підготовлено другий набір фрагментів зі стороною $S = 520$~пікселів, що дорівнює ширині контуру огляду камери ($\approx 518$~\px{}) з округленням до кратного 8, і перекриттям $o = 0{,}33$:
\begin{equation}
s = \lfloor 520 \cdot 0{,}67 \rfloor = 348, \qquad
N = \left( \left\lceil \frac{6200 - 520}{348} \right\rceil + 1 \right) \cdot \left( \left\lceil \frac{4320 - 520}{348} \right\rceil + 1 \right) = 18 \cdot 12 = 216 .
\end{equation}
Фрагменти залишено квадратними: під час розворотів БпЛА контур огляду на карті змінює орієнтацію з горизонтальної ($518 \times 292$) на вертикальну ($292 \times 518$), і прямокутний фрагмент підходив би лише для одного напрямку польоту. Обидва набори фрагментів зберігаються окремо, кожен --- зі своєю таблицею метаданих і кешем ознак.

\textbf{Перевірка цілісності даних.}\\
Після розбиття автоматично перевіряється, що кожен фрагмент попіксельно збігається з відповідною ділянкою вихідної карти, усі фрагменти лежать у межах карти, їх кількість відповідає формулі, а таблиця метаданих не містить пропусків і дублікатів.

\subsection{Вилучення кадрів із відео}

З відео вибирається один кадр на секунду: кожен $n$-й кадр, де $n = \operatorname{round}(f_{\text{відео}} / f_{\text{ціль}}) = \operatorname{round}(29{,}97 / 1) = 30$. Часова мітка кадру обчислюється як $t = i / f_{\text{відео}}$, де $i$ --- номер кадру у відео. Кадри декодуються послідовно (без позиціювання за номером, яке ненадійне для кодеків H.264/HEVC) і зберігаються без втрат у форматі PNG. У результаті отримано 95 кадрів із кроком $\approx 1{,}001$~с.

\subsection{Нормалізація масштабу та єдина система координат}
\label{sec:scale}

Кадр і ортофотоплан мають різну просторову роздільність: один піксель кадру відповідає приблизно $s_f = 0{,}27$ пікселя карти (значення визначено калібруванням, див. п.~\ref{sec:exp-pair}). Якщо подати в метод зіставлення зображення з різним масштабом деталей, частина відповідностей втрачається. Тому перед зіставленням обидва зображення приводяться до спільної роздільності:
\begin{itemize}
    \item фрагмент карти масштабується з коефіцієнтом $k$ (у роботі $k = 1$, тобто фрагмент подається як $512 \times 512$);
    \item кадр масштабується з коефіцієнтом $s_f \cdot k$: кадр $1920 \times 1080$ перетворюється на $\approx 518 \times 292$ пікселі. Кадр використовується повністю, без обрізання.
\end{itemize}

Методи зіставлення повертають координати відповідностей у пікселях зменшених (<<робочих>>) зображень. Система переносить їх назад в оригінальні пікселі кадру та фрагмента з урахуванням узгодження центрів пікселів, яке використовується під час масштабування:
\begin{equation}
x_{\text{ориг}} = \frac{x_{\text{роб}} + 0{,}5}{s_x} - 0{,}5, \qquad
y_{\text{ориг}} = \frac{y_{\text{роб}} + 0{,}5}{s_y} - 0{,}5,
\label{eq:backmap}
\end{equation}
де $s_x, s_y$ --- фактичні коефіцієнти масштабування за осями (вони можуть дещо відрізнятися від заданого через округлення розмірів). Унаслідок цього гомографія $H$ оцінюється між оригінальними пікселями кадру та оригінальними пікселями фрагмента. Усі геометричні величини (поріг RANSAC, похибка репроєкції, координати центра й контуру кадру) виражаються в пікселях карти і не залежать ні від методу зіставлення, ні від робочого масштабу.

\subsection{Реалізація методів зіставлення}

\textbf{Спільний інтерфейс.}\\
Обидва методи реалізовано як класи зі спільним інтерфейсом, що складається з двох кроків:
\begin{enumerate}
    \item \code{extract(image)} --- підготовка одного зображення: для LightGlue --- виділення ключових точок і дескрипторів SuperPoint; для LoFTR, який не має окремого етапу детекції, --- лише формування тензора;
    \item \code{match\_features(f\textsubscript{0}, f\textsubscript{1})} --- встановлення відповідностей між двома підготовленими зображеннями.
\end{enumerate}
Таке розділення дає змогу обчислювати ознаки кадру один раз під час перебору багатьох фрагментів, а ознаки фрагментів карти --- заздалегідь. Решта системи працює лише з цим інтерфейсом і не залежить від конкретного методу.

\textbf{LightGlue + SuperPoint.}\\
SuperPoint виділяє до 2048 ключових точок на зображення; LightGlue використовується з параметрами авторів: поріг адаптивної зупинки 0,95, поріг відсікання точок 0,99, поріг прийняття відповідності 0,1. Під час аналізу прототипу курсової роботи виявлено, що функція виділення ознак SuperPoint за замовчуванням масштабує вхідне зображення до 1024~пікселів за довшою стороною. Унаслідок цього в курсовій роботі LightGlue фактично працював із зображеннями $1024 \times 1024$, а LoFTR --- з $512 \times 512$, тобто методи порівнювалися за різних умов. У розробленій системі це неявне масштабування вимкнено, і масштаб для обох методів задається однаково (п.~\ref{sec:scale}).

\textbf{LoFTR.}\\
Використано реалізацію бібліотеки Kornia з вагами \textit{outdoor}; зображення перетворюються у відтінки сірого. Грубий рівень LoFTR працює на сітці з кроком 8 пікселів, тому розміри вхідних зображень округлюються до кратних 8; відповідні фактичні коефіцієнти масштабу враховуються у формулі~(\ref{eq:backmap}).

\textbf{Кешування ознак фрагментів карти.}\\
Для LightGlue ознаки SuperPoint усіх 384 фрагментів обчислюються один раз і зберігаються на диску (координати точок, оцінки та дескриптори у форматі float32, без втрати точності): $523$--$1609$ точок на фрагмент, загальний обсяг $\approx 347$~МБ, час побудови $\approx 14$~с. Разом із кешем зберігаються параметри, з якими його побудовано (параметри SuperPoint, робочий масштаб, контрольні суми таблиці фрагментів і карти); кеш, що не відповідає поточним параметрам, система відмовляється використовувати. Для набору з 216 фрагментів $520 \times 520$ кеш містить $485$--$1540$ точок на фрагмент, має обсяг $\approx 202$~МБ і будується за $\approx 10$~с. Для LoFTR кешування неможливе, оскільки його ознаки залежать від обох зображень пари.

\textbf{Вимірювання часу.}\\
Перед вимірюваннями виконується <<прогрівання>> моделей на зображеннях тих самих розмірів; для врахування асинхронних обчислень на GPU кожен вимірюваний блок обрамлюється синхронізацією CUDA. Час розділено на етапи: обробка кадру, обробка фрагмента, зіставлення, геометрична перевірка. Читання зображень із диска до вимірюваного часу не входить.

\subsection{Геометрична перевірка відповідностей}

\textbf{Оцінювання гомографії.}\\
За отриманими відповідностями $\{(\mathbf{p}_i, \mathbf{q}_i)\}$ (точка кадру --- точка фрагмента) оцінюється гомографія $H$ (див. розділ~1). Використовується алгоритм RANSAC \cite{ransac} із порогом $\tau = 3$~\px{}, максимальною кількістю ітерацій 5000 і довірчою ймовірністю 0,999; для відтворюваності генератор випадкових чисел ініціалізується фіксованим значенням. Як альтернативний оцінювач реалізовано MAGSAC++ \cite{magsac}, який замість жорсткого порогу зважує відповідності за їхньою похибкою; порівняння оцінювачів наведено в розділі~\ref{sec:experiments}.

Відповідність вважається \emph{геометрично узгодженою}, якщо пряма похибка репроєкції не перевищує порогу:
\begin{equation}
e_i = \left\| \pi(H \tilde{\mathbf{p}}_i) - \mathbf{q}_i \right\|_2 \le \tau,
\end{equation}
де $\tilde{\mathbf{p}}_i$ --- однорідні координати точки, $\pi([u, v, w]^T) = (u/w,\, v/w)$ --- перехід від однорідних координат до декартових. Якщо кількість відповідностей менша за мінімально необхідну для успіху (15), RANSAC не запускається: на наборі, що складається лише з хибних пар, \code{cv2.RANSAC} відпрацьовує всі 5000 ітерацій ($\approx 0{,}1$~с) без шансу на успішний результат.

\textbf{Перевірка правдоподібності гомографії.}\\
RANSAC може знайти формально узгоджену, але фізично неможливу модель (наприклад, на повторюваних структурах). Тому для знайденої $H$ обчислюється якобіан $J$ відображення в центрі кадру $\mathbf{c}$ --- локальне афінне наближення гомографії --- і з нього визначаються:
\begin{align}
\sigma &= \sqrt{|\det J|} && \text{--- масштаб (пікселів карти на піксель кадру)}, \\
\theta &= \operatorname{atan2}(J_{21} - J_{12},\; J_{11} + J_{22}) && \text{--- кут повороту кадру відносно карти}, \\
\kappa &= \sigma_{\max}(J) / \sigma_{\min}(J) && \text{--- анізотропія (відношення сингулярних чисел)}.
\end{align}
Гомографія визнається правдоподібною, якщо виконуються всі умови:
\begin{enumerate}
    \item третя однорідна координата має однаковий знак в усіх кутах кадру (рамка кадру не перетинає <<лінію нескінченності>>);
    \item $\det J > 0$ --- немає дзеркального відображення;
    \item проєкція рамки кадру є опуклим чотирикутником;
    \item $\kappa \le 1{,}5$ --- для надирної камери масштаб за осями майже однаковий;
    \item $s_f / 2 \le \sigma \le 2 s_f$ --- масштаб узгоджується з очікуваним.
\end{enumerate}
Кут $\theta$ надалі використовується як оцінка курсу БпЛА відносно карти, а масштаб $\sigma$ --- як непрямий показник висоти польоту.

\textbf{Критерій успішного зіставлення пари.}\\
Зіставлення пари <<кадр --- фрагмент>> вважається успішним, якщо гомографію знайдено, кількість геометрично узгоджених відповідностей $N_{\text{in}} \ge 15$ і гомографія правдоподібна. Інакше система фіксує конкретну причину невдачі.

\subsection{Автоматичний пошук фрагмента карти}
\label{sec:search}

Для кадру, положення якого невідоме, конвеєр зіставлення виконується з кожним фрагментом пулу (усієї карти, заданої прямокутної області або прорідженої сітки фрагментів). Ознаки кадру обчислюються один раз. Вибір найкращого фрагмента виконується у два кроки.

\textbf{Крок 1: жорсткий фільтр.}
Фрагмент-кандидат відкидається, якщо не виконується хоча б одна умова:
\begin{enumerate}
    \item гомографію знайдено і вона правдоподібна;
    \item $N_{\text{in}} \ge 15$;
    \item частка геометрично узгоджених відповідностей $r = N_{\text{in}} / N_{\text{all}} \ge 0{,}15$;
    \item середня похибка репроєкції $\bar{e} \le 3{,}5$~\px{};
    \item проєкція центра кадру $(x_{\text{лок}}, y_{\text{лок}})$ потрапляє в межі фрагмента з допуском $m = 30$~\px{}: $-m \le x_{\text{лок}} \le S - 1 + m$ (аналогічно для $y$), тобто фрагмент є основним носієм центра кадру;
    \item проєкція центра кадру лежить у межах карти.
\end{enumerate}
Умова~4 за порогу RANSAC $\tau = 3$~\px{} виконується завжди і є запобіжником на випадок зміни порогу.

\textbf{Крок 2: бал ранжування.}
Серед кандидатів, що пройшли фільтр, обирається фрагмент із найбільшим значенням
\begin{equation}
\operatorname{Score} = N_{\text{in}} \cdot \frac{N_{\text{in}}}{N_{\text{all}}} \cdot \exp\!\left(-\frac{\bar{e}}{\sigma_r}\right), \qquad \sigma_r = 2~\text{\px{}}.
\label{eq:score}
\end{equation}
Бал зростає з кількістю та часткою геометрично узгоджених відповідностей і плавно зменшується з похибкою репроєкції. На відміну від відношення $N_{\text{in}} / \bar{e}$, він не стає необмежено великим за малої похибки на кількох випадкових точках.

\textbf{Штраф за відхилення масштабу (вдосконалення).}\\
Перевірка правдоподібності допускає відхилення масштабу гомографії від очікуваного вдвічі, тому фрагмент, гомографію якого оцінено за площиною даху високої будівлі, може пройти фільтр і перемогти за балом. Так, для кадру 2 переміг фрагмент 220 із масштабом $\sigma = 0{,}207$, тоді як сусідній фрагмент 196 із масштабом $\sigma = 0{,}267$, близьким до калібрувального, мав нижчий бал (11,6 проти 13,2). Оскільки висота польоту приблизно стала, очікуваний масштаб $s_f$ відомий, і бал можна доповнити необов'язковим множником
\begin{equation}
\operatorname{Score}' = \operatorname{Score} \cdot \exp\!\left( -\frac{\ln^2(\sigma / s_f)}{2 \sigma_s^2} \right),
\label{eq:scale-prior}
\end{equation}
де $\sigma_s$ --- допустиме відносне відхилення масштабу (у роботі $\sigma_s = 0{,}15$, тобто $\approx \pm 15\,\%$). Логарифм робить штраф симетричним: відхилення в $0{,}8$ і $1{,}25$ раза штрафуються однаково. Для кадру 2 множник становить 0,21 для фрагмента 220 і $\approx 1{,}00$ для фрагмента 196, тож переможцем стає фрагмент, гомографія якого описує площину землі. Штраф м'який і не відкидає кандидатів, тому за помірної зміни висоти польоту правильний фрагмент лише втрачає частину балу. В основній конфігурації системи штраф вимкнено; його вплив досліджено в п.~\ref{sec:exp-improve}.

\textbf{Обчислення координат на карті.}\\
Центр кадру $\mathbf{c} = ((W_f - 1)/2,\, (H_f - 1)/2)$ (у конвенції центрів пікселів) переноситься гомографією обраного фрагмента з явною дегомогенізацією:
\begin{equation}
\tilde{\mathbf{p}} = H \begin{bmatrix} c_x \\ c_y \\ 1 \end{bmatrix}, \qquad
x_{\text{лок}} = \frac{\tilde{p}_1}{\tilde{p}_3}, \quad y_{\text{лок}} = \frac{\tilde{p}_2}{\tilde{p}_3},
\end{equation}
після чого обчислюються глобальні координати на карті:
\begin{equation}
X = X_0 + x_{\text{лок}}, \qquad Y = Y_0 + y_{\text{лок}},
\end{equation}
де $(X_0, Y_0)$ --- координати кута фрагмента з таблиці метаданих. Аналогічно переносяться чотири кути кадру, що дає контур огляду камери на карті. Точки, для яких $|\tilde{p}_3| \approx 0$, вважаються невизначеними й відкидаються.

\subsection{Локалізація послідовності кадрів}
\label{sec:tracking}

Повний перебір 384 фрагментів для кожного кадру є надлишковим: за 1~с апарат зміщується лише на кілька десятків пікселів карти. Тому для послідовності кадрів реалізовано трекер зі станами \textsc{lost} та \textsc{tracking}.

\textbf{Стан \textsc{lost}} (перший кадр або втрата супроводу): виконується глобальний пошук (п.~\ref{sec:search}). Якщо переможець має бал $\operatorname{Score} \ge 5$, положення приймається (статус \textsc{initialized}) і трекер переходить у стан \textsc{tracking}.

\textbf{Стан \textsc{tracking}}: пошук виконується лише серед фрагментів вікна $3 \times 3$ сітки (не більше 9 фрагментів) навколо базового фрагмента. Базовим обирається фрагмент, центр якого найближчий до \emph{прогнозованого} положення. Прогноз будується за двома останніми прийнятими положеннями $(t_{k-2}, P_{k-2})$, $(t_{k-1}, P_{k-1})$ у припущенні сталої швидкості:
\begin{equation}
\hat{P}_k = P_{k-1} + \mathbf{v} \cdot (t_k - t_{k-1}), \qquad
\mathbf{v} = \frac{P_{k-1} - P_{k-2}}{t_{k-1} - t_{k-2}},
\label{eq:extrapolation}
\end{equation}
де модуль швидкості обмежується значенням $V_{\max}$. За сталого інтервалу між кадрами це збігається з $\hat{P}_k = P_{k-1} + (P_{k-1} - P_{k-2})$.

\textbf{Кінематична перевірка.}\\
Нехай $P_{\text{ост}}$ --- останнє \emph{прийняте} положення, отримане в момент $t_{\text{ост}}$. Нова оцінка $P_k$ приймається, якщо
\begin{equation}
\| P_k - P_{\text{ост}} \|_2 \le D_{\max} = V_{\max} \cdot (t_k - t_{\text{ост}}) + D_{\text{запас}},
\end{equation}
де $V_{\max} = 120$~\px{}/с, $D_{\text{запас}} = 30$~\px{}; за частоти 1~кадр/с $D_{\max} = 150$~\px{}. Порівняння саме з останнім прийнятим положенням означає, що після кількох пропущених кадрів допустимий радіус зростає пропорційно часу. Серед кандидатів, що пройшли фільтр п.~\ref{sec:search}, обирається найкращий за балом, який проходить і кінематичну перевірку. Якщо такого кандидата немає, положення кадру екстраполюється за формулою~(\ref{eq:extrapolation}) (статус \textsc{interpolated}), а відкинута оцінка зберігається для аналізу. Після трьох поспіль кадрів без прийнятого положення трекер переходить у стан \textsc{lost}.

\textbf{Компенсація курсу.}\\
SuperPoint/LightGlue і LoFTR не інваріантні до великих поворотів зображення (кількісне підтвердження наведено в п.~\ref{sec:exp-robust}), а під час розвороту БпЛА кадр повертається відносно карти. Кут <<кадр $\to$ карта>> відомий із гомографій прийнятих кадрів ($\theta$, п.~\ref{sec:search}), тому перед зіставленням кадр повертається на прогнозований кут $\hat{\theta}_k$, і метод зіставлення отримує кадр, орієнтований так само, як карта. Прогноз кута будується аналогічно до~(\ref{eq:extrapolation}) з урахуванням переходу через $\pm 180^\circ$, а кутова швидкість обмежується значенням $45^\circ$/с. Полотно під час повороту розширюється, а відповідності, що потрапили на фон поза кадром, відкидаються. Перетворення <<оригінальний кадр $\to$ робоче зображення>> є композицією масштабування та повороту $T = R(\hat{\theta}_k) \cdot S$, і координати відповідностей переносяться назад точним оберненим перетворенням $T^{-1}$. Тому гомографія, як і раніше, оцінюється в оригінальних пікселях кадру.

Якщо з прогнозованим курсом фрагмент не знайдено, у стані \textsc{tracking} послідовно перевіряються гіпотези $\hat{\theta}_k \pm 45^\circ$ і $\hat{\theta}_k \pm 90^\circ$. У стані \textsc{lost}, де повний перебір карти дорогий, кожен наступний глобальний пошук використовує чергову гіпотезу з циклу $0^\circ, 90^\circ, 180^\circ, -90^\circ$ відносно останнього відомого курсу.

\textbf{Вікно карти за прогнозом (вдосконалення).}\\
Пошук у вікні $3 \times 3$ має два недоліки. По-перше, кадр зіставляється з кожним із 9 фрагментів окремо, хоча положення апарата вже приблизно відоме. По-друге, контур огляду ($\approx 518 \times 292$~\px{}) не вміщується у фрагмент, тому кожен фрагмент містить лише частину кадру: відповідності знаходяться на одній стороні кадру, а його центр не <<оточений>> точками, що погіршує обумовленість гомографії. Межі фрагментів задано сіткою, а не положенням апарата, тому результат залежить також від того, як кадр ліг на сітку.

У стані \textsc{tracking} відомі прогнозовані положення $\hat{P}_k$ і курс $\hat{\theta}_k$, а форма контуру огляду між сусідніми кадрами майже не змінюється. Тому кадр спершу зіставляється з однією ділянкою карти, вирізаною навколо очікуваного контуру (рис.~\ref{fig:window}):
\begin{enumerate}
    \item контур останнього прийнятого кадру $Q_{\text{ост}}$ (чотири вершини на карті) переноситься в прогнозоване положення: $\hat{Q}_k = Q_{\text{ост}} - P_{\text{ост}} + \hat{P}_k$;
    \item вікном є охопний прямокутник $\hat{Q}_k$, розширений на запас $m_w = 128$~\px{} з кожного боку й обрізаний межами карти; сторони вікна округлюються до кратних 8 (вимога LoFTR), а вікно зі стороною понад 1024~\px{} не використовується;
    \item кадр, повернутий на прогнозований курс $\hat{\theta}_k$, зіставляється з вікном тим самим конвеєром; до кандидата застосовуються ті самі фільтр (п.~\ref{sec:search}) і кінематична перевірка;
    \item якщо оцінку не прийнято, виконується звичайний пошук у вікні $3 \times 3$ фрагментів з усіма гіпотезами курсу.
\end{enumerate}
Запас $m_w$ покриває похибку прогнозу положення (крок польоту в середньому 44~\px{} за кадр) і зміну охопного прямокутника контуру під час повороту. На досліджуваному польоті курс між сусідніми кадрами змінюється з медіаною $0{,}5^\circ$ і 90-м процентилем $14^\circ$; за повороту на $20^\circ$ висота охопного прямокутника контуру $518 \times 292$ зростає на $\approx 160$~\px{}, тобто на $\approx 80$~\px{} з кожного боку, що вкладається в запас. На різких розворотах (до $41^\circ$ між кадрами) кадр може не вміститися у вікно, і тоді спрацьовує резервний пошук у фрагментах. Вікно має розмір $\approx 770 \times 550$~пікселів, і за звичайного польоту кадр потрапляє в нього повністю. Ознаки вікна обчислюються під час обробки кадру (кешування неможливе, оскільки вікно щоразу інше), однак замість до 9 зіставлень виконується одне.

\begin{figure}[H]
\centering
\begin{tikzpicture}[scale=0.0062, line join=round]
  % ліворуч: вікно 3x3 фрагментів 512 з кроком 256
  \begin{scope}
    \foreach \i in {0,1,2} \foreach \j in {0,1,2}
      \draw[gray!55] (\i*256, \j*256) rectangle ++(512, 512);
    \draw[blue!70!black, very thick] (256, 256) rectangle ++(512, 512);
    \fill[green!45!black, opacity=0.25] (380, 470) rectangle ++(518, 292);
    \draw[green!45!black, thick] (380, 470) rectangle ++(518, 292);
    \draw[red, thick] (619, 596) -- ++(40, 40) (619, 636) -- ++(40, -40);
    \node[font=\small, align=center] at (512, -110) {вікно $3\times3$: до 9 зіставлень,\\кадр частково поза фрагментом};
  \end{scope}
  % праворуч: одне вікно навколо прогнозованого контуру
  \begin{scope}[xshift=1450]
    \draw[gray!55] (0, 0) rectangle (1024, 1024);
    \draw[orange!85!black, very thick] (252, 342) rectangle ++(774, 548);
    \fill[green!45!black, opacity=0.25] (380, 470) rectangle ++(518, 292);
    \draw[green!45!black, thick] (380, 470) rectangle ++(518, 292);
    \draw[red, thick] (619, 596) -- ++(40, 40) (619, 636) -- ++(40, -40);
    \draw[<->, thin] (252, 920) -- node[above, font=\scriptsize] {$m_w$} (380, 920);
    \node[font=\small, align=center] at (512, -110) {вікно за прогнозом: 1 зіставлення,\\кадр повністю у вікні};
  \end{scope}
\end{tikzpicture}
\caption{Пошук у стані \textsc{tracking}: синім --- один із фрагментів сітки, помаранчевим --- вікно карти навколо прогнозованого контуру огляду (зелений), червоним --- прогнозоване положення}
\label{fig:window}
\end{figure}

\textbf{Уточнення положення після глобального пошуку (вдосконалення).}\\
Глобальний пошук також виконується на фрагментах сітки, тому кадр-переможець, як правило, обрізаний краєм фрагмента. Для уточнення переможця кадр повторно зіставляється з ділянкою карти, що охоплює його контур огляду із запасом 64~\px{}. Уточнена оцінка приймається, якщо вона проходить фільтр п.~\ref{sec:search}, має не менше геометрично узгоджених відповідностей, ніж переможець, і зміщує центр кадру не більше ніж на 30~\px{} --- тобто є уточненням, а не новою гіпотезою положення. Інакше зберігається результат глобального пошуку.

\textbf{Курс із компаса (вдосконалення).}\\
Компенсація курсу за гомографією працює лише після першої успішної прив'язки. У стані \textsc{lost} курс невідомий, і глобальний пошук перебирає гіпотези $0^\circ, 90^\circ, 180^\circ, -90^\circ$, тобто курс кадру може відрізнятися від використаного до $45^\circ$. Магнітометр (компас) є стандартним датчиком БпЛА, працює без супутникової навігації і дає курс у реальному часі. Тому в систему додано необов'язкове використання його показу. Кут <<кадр $\to$ карта>> за компасом обчислюється як
\begin{equation}
\theta_{\text{к}} = \psi + \Delta,
\label{eq:compass}
\end{equation}
де $\psi$ --- показ компаса (курс камери відносно півночі), $\Delta$ --- стала поправка між напрямком, відносно якого вимірюється курс, і віссю $y$ карти. Поправка складається з магнітного схилення (якщо компас дає магнітний, а не географічний курс) і кута зближення меридіанів картографічної проєкції. Для проєкції UTM цей кут дорівнює
\begin{equation}
\gamma = \arctan\bigl(\tan(\lambda - \lambda_0) \sin\varphi\bigr),
\end{equation}
де $\varphi, \lambda$ --- широта й довгота району польоту, $\lambda_0$ --- осьовий меридіан зони. Поправку обчислюють один раз до польоту, без використання еталона. Якщо камера закріплена з поворотом відносно корпусу, а компас вимірює курс корпусу, до $\Delta$ додається ще й кут встановлення камери.

Показ компаса використовується двічі:
\begin{itemize}
    \item у стані \textsc{lost} глобальний пошук виконується одразу з курсом $\theta_{\text{к}}$ замість чергової гіпотези циклу;
    \item у стані \textsc{tracking} першою, як і раніше, перевіряється гіпотеза $\hat{\theta}_k$ за гомографією, оскільки вона точніша за компас. Якщо $|\theta_{\text{к}} - \hat{\theta}_k| > 10^\circ$ (наприклад, після різкого розвороту прогноз застарів), курс компаса перевіряється другою гіпотезою, до гіпотез $\pm 45^\circ$ і $\pm 90^\circ$.
\end{itemize}
Остаточні положення й кут завжди визначаються зіставленням із картою, тому похибка компаса не переходить у похибку координат. Вона лише має не перевищувати допуску методу зіставлення до повороту ($\approx 20$--$30^\circ$, п.~\ref{sec:exp-robust}). Для перевірки стійкості до неточного компаса до показу можна додати стале зміщення та гауссів шум; шум генерується з фіксованим зерном і для кожного кадру однаковий у всіх прогонах.

Ці вдосконалення ввімкнюються параметрами конфігурації; в основній конфігурації, результати якої наведено в пп.~\ref{sec:exp-yaw}--\ref{sec:exp-robust}, їх вимкнено.

Узагальнено обробка одного кадру в стані \textsc{tracking} складається з таких кроків:
\begin{enumerate}
    \item обчислити прогноз положення $\hat{P}_k$ і курсу $\hat{\theta}_k$;
    \item якщо ввімкнено вікно за прогнозом --- зіставити кадр, повернутий на $\hat{\theta}_k$, з вікном карти навколо $\hat{Q}_k$; якщо кандидат проходить фільтр і кінематичну перевірку --- прийняти положення і завершити обробку кадру;
    \item сформувати пул із не більше ніж 9 фрагментів навколо фрагмента, найближчого до $\hat{P}_k$;
    \item для кожної гіпотези курсу, починаючи з $\hat{\theta}_k$ (далі --- курс компаса, якщо він заданий і відрізняється від $\hat{\theta}_k$ більше ніж на $10^\circ$, потім $\hat{\theta}_k \pm 45^\circ$, $\hat{\theta}_k \pm 90^\circ$):
    \begin{enumerate}
        \item повернути та масштабувати кадр, обчислити його ознаки;
        \item зіставити кадр із фрагментами пулу, виконати геометричну перевірку, фільтр і ранжування;
        \item якщо найкращий кандидат проходить кінематичну перевірку --- прийняти положення (статус \textsc{tracking}) і завершити обробку кадру;
    \end{enumerate}
    \item якщо жодна гіпотеза не дала прийнятного положення --- екстраполювати положення (статус \textsc{interpolated}); після трьох таких кадрів поспіль перейти в стан \textsc{lost}.
\end{enumerate}

\subsection{Згладжування траєкторії фільтром Калмана}
\label{sec:kalman}

Трекер оцінює положення кожного кадру незалежно, тому траєкторія містить випадкові відхилення від кадру до кадру. Для їх зменшення реалізовано згладжування траєкторії фільтром Калмана \cite{kalman} з моделлю руху зі сталою швидкістю. Згладжування застосовується до готової траєкторії і не впливає на рішення трекера, що дає змогу порівняти згладжену та незгладжену траєкторії на тих самих прогонах.

Стан фільтра $\mathbf{x}_k = [X, Y, v_X, v_Y]^T$ містить положення на карті (\px{}) і швидкість (\px{}/с). Між кадрами з інтервалом $\Delta t$ стан змінюється за моделлю
\begin{equation}
\mathbf{x}_k = F \mathbf{x}_{k-1} + \mathbf{w}_k, \qquad
F = \begin{bmatrix} I & \Delta t \, I \\ 0 & I \end{bmatrix}, \qquad
Q = \sigma_a^2 \begin{bmatrix} \tfrac{\Delta t^4}{4} I & \tfrac{\Delta t^3}{2} I \\ \tfrac{\Delta t^3}{2} I & \Delta t^2 I \end{bmatrix},
\end{equation}
де $I$ --- одинична матриця $2 \times 2$, $\mathbf{w}_k$ --- шум процесу (дискретний білий шум прискорення з середньоквадратичним відхиленням $\sigma_a$), $Q$ --- його коваріаційна матриця. Вимірюванням є положення, знайдене трекером:
\begin{equation}
\mathbf{z}_k = [X_k, Y_k]^T = [\, I \;\; 0 \,] \, \mathbf{x}_k + \mathbf{v}_k, \qquad R = \sigma_z^2 I .
\end{equation}
Вимірюваннями вважаються лише кадри зі статусами \textsc{tracking} та \textsc{initialized}. Для кадрів \textsc{interpolated} (екстраполяція трекера) використовується прогноз фільтра, кадри \textsc{lost} залишаються без положення, а кожна повторна ініціалізація розпочинає новий сегмент траєкторії з новою ініціалізацією фільтра.

Параметри задано апріорно, а не підібрано за еталоном, оскільки підбір за тими самими опорними кадрами завищив би оцінку ефекту: $\sigma_z = 4$~\px{} відповідає типовій похибці одиничної оцінки, а $\sigma_a = 20$~\px{}/с$^2$ --- маневру з розворотом на $90^\circ$ за $\approx 3$~с за швидкості $\approx 44$~\px{}/с. Обчислюються дві оцінки:
\begin{itemize}
    \item \emph{фільтр Калмана} використовує лише попередні кадри і може працювати в реальному часі;
    \item \emph{згладжувач Рауха--Тунга--Штрібеля} (RTS) \cite{rts} додатково виконує зворотний прохід по сегменту, тобто використовує й наступні кадри, і придатний лише для аналізу польоту після його завершення.
\end{itemize}

\subsection{Набори даних і географічна прив'язка}
\label{sec:datasets}

\textbf{Набори даних.}\\
Щоб перевіряти систему на різних польотах без зміни коду, вхідні дані оформлено як набори даних. Кожен набір описується файлом, у якому вказано ортофотоплан, відео, (необов'язково) файл географічної прив'язки та журнал польоту, а також параметри, що залежать від даних: масштаб кадру $s_f$, розмір і перекриття фрагментів, розмір пікселя карти, поправку компаса. Набір обирається змінною середовища. Кадри, фрагменти, кеш ознак, еталон і результати кожного набору зберігаються в окремих каталогах, тому дані й результати першого польоту під час роботи з другим не змінюються. Алгоритм локалізації для всіх наборів однаковий.

\textbf{Географічна прив'язка карти.}\\
Прив'язку ортофотоплану задано файлом прив'язки (\textit{world file}), який містить розмір пікселя $a$, $e$ (для карти з північчю вгорі $e = -a$) і координати $(c, f)$ центра лівого верхнього пікселя в картографічній проєкції. Перехід від координат проєкції $(X_{\text{UTM}}, Y_{\text{UTM}})$ до пікселів карти:
\begin{equation}
x = \frac{X_{\text{UTM}} - c}{a}, \qquad y = \frac{Y_{\text{UTM}} - f}{e} .
\end{equation}
Прив'язка не використовується під час локалізації. Вона потрібна лише для двох допоміжних задач: перенесення еталонних координат GNSS на карту та переведення похибки з пікселів у метри ($E_{\text{м}} = a \cdot E_{\text{\px{}}}$).

\textbf{Еталон за журналом польоту.}\\
Для другого польоту еталонне положення кожного кадру отримано з журналу польоту DJI, у якому 10 разів на секунду записано широту, довготу, висоту й орієнтацію апарата. Для кадру з часовою міткою $t$ береться запис журналу з часом польоту $t + t_0$, де $t_0$ --- момент початку відеозапису. Широта й довгота переводяться в координати UTM (за рядами Крюгера), а потім у пікселі карти за файлом прив'язки. Журнал використовується лише для оцінювання вже отриманої траєкторії: система локалізації його не читає. Похибка самого GNSS-положення становить $\approx 2$--$3$~м, тому такий еталон менш точний за ручну розмітку першого польоту, але доступний для кожного кадру.

\textbf{Показ компаса.}\\
Для експериментів із компасом з журналу береться \emph{лише} курс камери (координати не використовуються). Курс інтерполюється на моменти кадрів із розгортанням кута через $360^\circ$ і зберігається окремим файлом, який імітує бортовий датчик. Поправку $\Delta$, зміщення та шум (формула~\ref{eq:compass}) застосовує трекер.

\subsection{Модульна структура та відтворюваність}

Програмну систему організовано як пакет із незалежними модулями (табл.~\ref{tab:modules}) і набір сценаріїв командного рядка для кожного етапу.

\begin{table}[H]
\centering
\caption{Основні модулі програмної системи}
\label{tab:modules}
\small
\begin{tabular}{p{0.36\linewidth} p{0.56\linewidth}}
\toprule
\textbf{Модуль} & \textbf{Призначення} \\
\midrule
\code{preprocessing/tiles.py}, \code{frames.py} & розбиття карти, вилучення кадрів, метадані \\
\code{preprocessing/images.py} & читання/запис, масштабування, поворот, перенесення координат, спотворення для експериментів \\
\code{preprocessing/cache.py} & дисковий кеш ознак фрагментів \\
\code{matching/lightglue.py}, \code{loftr.py} & реалізації методів зі спільним інтерфейсом \\
\code{matching/pipeline.py} & єдиний конвеєр <<кадр + фрагмент $\to$ відповідності $\to$ RANSAC $\to$ метрики>> \\
\code{geometry/homography.py} & RANSAC/MAGSAC++, похибка репроєкції, перевірка правдоподібності \\
\code{evaluation/scoring.py} & фільтр і бал ранжування фрагментів, штраф за відхилення масштабу \\
\code{localization/coordinates.py} & перенесення точок на карту, контур огляду \\
\code{localization/search.py}, \code{tracker.py} & пошук серед фрагментів і довільних ділянок карти, трекер послідовності, вікно за прогнозом, уточнення положення \\
\code{localization/smoothing.py} & фільтр Калмана і згладжувач RTS для траєкторії \\
\code{localization/compass.py} & показ компаса, поправка $\Delta$, імітація зміщення й шуму \\
\code{evaluation/metrics.py}, \code{ground\_truth.py} & метрики, еталонні координати \\
\code{evaluation/geo.py} & перетворення в UTM, файл прив'язки карти, розмір пікселя \\
\code{visualization/} & рисунки відповідностей, положення та траєкторій \\
\bottomrule
\end{tabular}
\end{table}

Для відтворюваності результатів кожен запуск записує, разом із результатами, файл метаданих з усіма параметрами конвеєра й методу зіставлення, ідентифікатором версії коду та відомостями про обчислювальне середовище. Генератори випадкових чисел RANSAC і спотворень ініціалізуються фіксованими значеннями. Коректність перенесення координат, оцінювання гомографії, критеріїв фільтрації, логіки трекера та вдосконалень (штрафу за масштаб, вікна за прогнозом, уточнення положення, фільтра Калмана, курсу з компаса), а також перетворень координат UTM і файлу прив'язки перевіряється набором із 87 автоматичних тестів на синтетичних даних, для яких точний результат відомий заздалегідь; тести не потребують нейромереж і виконуються за кілька секунд.

\textbf{Серії експериментів.}\\
Щоб порівнювати варіанти системи, не змінюючи основних результатів, реалізовано окремі серії експериментів. Набір фрагментів (разом із таблицею метаданих і кешем ознак) і каталог результатів обираються змінними середовища; прогін на іншому наборі фрагментів завжди записується в окремий каталог. Сценарій серії запускає трекінг для кожної з чотирьох конфігурацій у кожному з варіантів системи --- базовому, з окремим вдосконаленням або з усіма вдосконаленнями разом. Кожен варіант задає всі перемикачі явно, тому результат не залежить від значень за замовчуванням у конфігурації. Зведені таблиці будуються для всіх варіантів, а для кожного з них наводиться різниця з базовим варіантом.

\subsection{Візуалізація результатів}

Система формує такі графічні матеріали:
\begin{itemize}
    \item відповідності між кадром і фрагментом карти: геометрично узгоджені --- зеленим, решта --- червоним;
    \item контур огляду камери і центр кадру на ортофотоплані разом із сіткою перевірених фрагментів і обраним фрагментом;
    \item траєкторію польоту на карті з позначенням статусів кадрів і напрямку руху;
    \item графіки похибок, часу обробки та результатів експериментів (розділ~\ref{sec:experiments}).
\end{itemize}





% =============================================================================
% Розділ 3. Результати експериментів та їх оцінка
% Підключення: \input{sections/section3_experiments}
% =============================================================================
\providecommand{\px}{пкс}
\providecommand{\code}[1]{\texttt{#1}}

\newpage
\section{Результати експериментів та їх оцінка}
\label{sec:experiments}

\subsection{Методика експериментального дослідження}

\textbf{Дані та обладнання.}\\
Усі експерименти виконано на даних, описаних у розділі~\ref{sec:implementation}: ортофотоплан $6200 \times 4320$~пікселів, розбитий на 384 фрагменти $512 \times 512$ з перекриттям 50\,\%, і 95 кадрів відео польоту БпЛА з частотою 1~кадр/с. Обчислення виконано на GPU NVIDIA GeForce GTX 1650. Якщо не зазначено інше, використано параметри за замовчуванням: $s_f = 0{,}27$, $k = 1$, поріг RANSAC 3~\px{}, мінімум 15 геометрично узгоджених відповідностей. Пункти~\ref{sec:exp-pair}--\ref{sec:exp-terrain} описують основну конфігурацію системи без вдосконалень; вдосконалення досліджено окремою серією експериментів на альтернативному розбитті карти (п.~\ref{sec:exp-improve}). Перевірку системи на другому польоті над іншим типом місцевості описано в п.~\ref{sec:exp-flight2}.

\textbf{Досліджувані конфігурації.}\\
Порівнювалися чотири конфігурації системи --- поєднання методу зіставлення (LightGlue або LoFTR) та оцінювача гомографії (RANSAC або MAGSAC++). Решта етапів (масштабування, геометрична перевірка, фільтр і ранжування, трекер, компенсація курсу) для всіх конфігурацій однакові. Це забезпечує коректність порівняння: на відміну від курсової роботи, обидва методи отримують зображення однакового масштабу.

\textbf{Метрики зіставлення.}\\
Для пари <<кадр --- фрагмент>> використовуються:
\begin{itemize}
    \item $N_{\text{all}}$ --- кількість відповідностей, які повернув метод;
    \item $N_{\text{in}}$ --- кількість геометрично узгоджених відповідностей (похибка репроєкції не більша за поріг RANSAC);
    \item $r = N_{\text{in}} / N_{\text{all}}$ --- частка геометрично узгоджених відповідностей;
    \item $\bar{e}$ --- середня похибка репроєкції геометрично узгоджених відповідностей, \px{}. За побудовою вона не перевищує порогу RANSAC, тому характеризує точність відповідностей, а не успішність локалізації;
    \item час обробки, с.
\end{itemize}

\textbf{Метрики локалізації.}\\
Похибка локалізації кадру $i$ визначається відносно еталонного положення (п.~\ref{sec:gt}):
\begin{equation}
E_i = \sqrt{(X_i - X_i^{\text{ет}})^2 + (Y_i - Y_i^{\text{ет}})^2}, \quad \text{\px{}}.
\end{equation}
За сукупністю кадрів обчислюються середня, медіанна та максимальна похибка, середньоквадратична похибка $\text{RMSE} = \sqrt{\tfrac{1}{N}\sum_i E_i^2}$, 90-й процентиль $P_{90}$, а також частка успішних локалізацій
\begin{equation}
\text{Success@}K = \frac{\left|\{ i : E_i \le K \}\right|}{N}, \qquad K \in \{10, 25, 50\}~\text{\px{}},
\end{equation}
де $N$ --- кількість усіх кадрів з еталоном; кадри, для яких система не видала положення, вважаються неуспішними. Ортофотоплан не має просторової прив'язки, тому похибки наведено в пікселях карти і в метри не переводяться.

\subsection{Еталонні координати}
\label{sec:gt}

Для обчислення похибки локалізації потрібне еталонне положення кадру, незалежне від системи, що оцінюється. Використання як еталона результатів самої системи або їхнього візуального <<уточнення>> призвело б до заниженої оцінки похибки, тому еталон побудовано за ручними відповідностями точок.

\textbf{Процедура розмітки.}\\
Для опорного кадру оператор у спеціально розробленому інструменті вказує щонайменше 5 пар відповідних точок <<кадр $\leftrightarrow$ карта>>, розташованих на рівні землі (кути дорожньої розмітки та бордюрів, перехрестя доріжок, основи стовпів). Дахи будівель не використовуються, оскільки через паралакс вони зміщені відносно земної поверхні. За всіма парами методом найменших квадратів (без RANSAC) оцінюється гомографія $H_{\text{ет}}$, і центр кадру переноситься на карту. Приклад розмітки наведено на рисунку~\ref{fig:gt-example}.

\begin{figure}[H]
    \centering
    \includegraphics[width=\linewidth]{images/master/gt_annotation_example.png}
    \caption{Розмітка еталонного положення кадру 52 (кадр повернутий на $\approx 170^\circ$ відносно карти): пари точок пронумеровано, зеленим --- спроєктований контур кадру, червоним хрестиком --- центр кадру}
    \label{fig:gt-example}
\end{figure}

\textbf{Контроль якості розмітки.}\\
З чотирьох пар гомографія визначається точно, тому залишкова похибка дорівнює нулю й не дає змоги виявити помилку оператора. У першій версії розмітки саме такі кадри (по 4 пари, скупчені в одній частині кадру) давали <<похибку>> системи понад 100~\px{} для всіх методів одночасно. Тому опорний кадр вважається надійним, якщо:
\begin{enumerate}
    \item вказано не менше 5 пар;
    \item залишкова похибка кожної пари не перевищує 3~\px{}, інакше одна з точок хибна;
    \item 95-й процентиль зміщення центру за випадкового шуму кліків ($\sigma = 2$~пікселі кадру та 1~піксель карти, метод Монте-Карло) не перевищує 5~\px{}, інакше точки скупчені або центр кадру розташований далеко від них.
\end{enumerate}
Критерій не використовує результатів системи. Опорні кадри обрано візуально так, щоб на них були чіткі наземні орієнтири і щоб вони рівномірно покривали політ, включно з розворотами. Кадри над пустирем (33--40) і каналом (66--69) пропущено через відсутність надійних точок.

Розмічено 19 опорних кадрів, з них 18 надійних: 5--8 пар на кадр (у середньому 6,2), медіанна невизначеність центру 2,6~\px{} (максимум 5,0~\px{}), найбільша залишкова похибка пари --- 3,0~\px{}. Ненадійний кадр 40 (пустир, 4 пари) в метрики не включено. Для 7 надійних кадрів центр кадру лежить поза областю вказаних точок, тобто його положення екстрапольовано. Шум кліків цим критерієм враховано, але систематичні похибки (дисторсія ширококутного об'єктива, паралакс) за межами точок не оцінюються, що слід мати на увазі під час аналізу окремих кадрів.

Для кадрів між опорними еталон можна лінійно інтерполювати в часі, однак на розворотах інтерполяція <<зрізає>> траєкторію. Через це похибка за інтерпольованим еталоном вимірює насамперед похибку інтерполяції: для всіх чотирьох конфігурацій вона однакова, $\approx 38$~\px{}. Тому всі наведені нижче похибки обчислено лише за 18 надійними опорними кадрами.

\subsection{Зіставлення кадру з відомим фрагментом карти}
\label{sec:exp-pair}

\textbf{Результати для еталонної пари.}\\
Для перевірки конвеєра використано пару з курсової роботи: кадр 2 і фрагмент карти, що містить місце зйомки (результати --- у табл.~\ref{tab:pair}, приклади відповідностей --- на рис.~\ref{fig:pair}).

\begin{table}[H]
\centering
\caption{Зіставлення еталонної пари <<кадр 2 --- фрагмент 196>> (GPU, медіана часу з 5 запусків)}
\label{tab:pair}
\begin{tabular}{lcccccc}
\toprule
\textbf{Метод} & $N_{\text{all}}$ & $N_{\text{in}}$ & $r$ & $\bar{e}$, \px{} & масштаб $\sigma$ & час, с \\
\midrule
LightGlue & 310 & 89 & 0,29 & 1,58 & 0,267 & 0,25 \\
LoFTR & 631 & 184 & 0,29 & 1,40 & 0,202 & 0,27 \\
\bottomrule
\end{tabular}
\end{table}

\begin{figure}[H]
    \centering
    \begin{subfigure}{\linewidth}
        \centering
        \includegraphics[width=0.92\linewidth]{images/master/pair_lightglue_matches.png}
        \caption{LightGlue}
    \end{subfigure}
    \vspace{2mm}
    \begin{subfigure}{\linewidth}
        \centering
        \includegraphics[width=0.92\linewidth]{images/master/pair_loftr_matches.png}
        \caption{LoFTR}
    \end{subfigure}
    \caption{Відповідності між кадром 2 (ліворуч) і фрагментом карти (праворуч): зеленим --- геометрично узгоджені, червоним --- відкинуті RANSAC}
    \label{fig:pair}
\end{figure}

\textbf{Виправлення методики курсової роботи.}\\
У курсовій роботі на цій парі LightGlue мав 229 геометрично узгоджених відповідностей проти 117 у LoFTR і працював повільніше (6,55~с проти 4,18~с). Аналіз прототипу показав, що ці результати отримано за нерівних умов. По-перше, SuperPoint неявно масштабував вхідні зображення до 1024~пікселів, тоді як LoFTR працював із 512. По-друге, обидва методи бачили лише центральну квадратну частину кадру, без нормалізації масштабу. По-третє, час вимірювався без прогрівання моделей. За однакових умов (табл.~\ref{tab:pair}) співвідношення змінюється: LoFTR знаходить удвічі більше геометрично узгоджених відповідностей за однакової їх частки, а час обробки пари на GPU для обох методів становить $\approx 0{,}25$~с. Тому числові результати курсової роботи надалі не використовуються.

\textbf{Калібрування масштабу кадру.}\\
Початкову оцінку $s_f \approx 0{,}25$ отримано вручну. Для її уточнення кадр 2 зіставлено з фрагментом карти $1024 \times 1024$, що повністю містить контур кадру. Обидва методи за обох оцінювачів і робочих масштабів $k \in \{1;\, 1{,}5\}$ дали узгоджений масштаб гомографії $\sigma = 0{,}267$--$0{,}277$, а оцінки центра кадру розійшлися не більше ніж на 16~\px{}. Отже, прийнято $s_f = 0{,}27$: кадр $1920 \times 1080$ покриває на карті ділянку $\approx 518 \times 292$~\px{}, ширшу за фрагмент 512~\px{}.

\textbf{Вплив непланарності сцени.}\\
На фрагменті 196 масштаб гомографії LoFTR ($\sigma = 0{,}20$) суттєво менший за калібрувальний (табл.~\ref{tab:pair}). Геометрично узгоджені відповідності LoFTR зосереджені на даху високої будівлі, ближчому до камери, ніж земля, тому RANSAC оцінює гомографію площини даху, а не площини землі (рис.~\ref{fig:pair}б). Центр кадру при цьому все одно визначається з похибкою $\approx 15$~\px{}, оскільки він сам розташований на цьому даху. Цей ефект є прямим наслідком припущення про плоску сцену.

\textbf{Контрольна перевірка на хибних фрагментах.}\\
Кадр 2 зіставлено з 33 фрагментами, серед яких є як сусідні до місця зйомки, так і віддалені. Фрагменти, що не перетинаються з кадром, дали щонайбільше 15 (LightGlue) і 14 (LoFTR) геометрично узгоджених відповідностей і всі були відкинуті: частина за порогом $N_{\text{in}} \ge 15$, частина --- перевіркою правдоподібності гомографії. Усі фрагменти, що перетинаються з кадром, дали від 18 до 240 таких відповідностей і пройшли перевірку. Запас між цими групами невеликий, тому для вибору фрагмента додатково використовуються частка відповідностей і умова потрапляння центра кадру у фрагмент (п.~\ref{sec:search}).

\subsection{Автоматичний пошук фрагмента для одиночного кадру}

Для кадру 2 виконано пошук без будь-якої інформації про місце зйомки (табл.~\ref{tab:search}, рис.~\ref{fig:global-loc}).

\begin{table}[H]
\centering
\caption{Автоматичний пошук фрагмента для кадру 2}
\label{tab:search}
\small
\begin{tabular}{lccccc}
\toprule
\textbf{Запуск} & \textbf{Перевірено} & \textbf{Пройшли фільтр} & \textbf{Переможець} & \textbf{Центр $(X, Y)$, \px{}} & \textbf{Час, с} \\
\midrule
LightGlue, уся карта, кеш & 384 & 5 & 220 & (1121,4;\ 2344,7) & 33 \\
LoFTR, обмежена область & 56 & 5 & 197 & (1108,6;\ 2327,5) & 18 \\
LoFTR, кожен 2-й рядок і стовпець & 96 & 1 & 197 & (1108,6;\ 2327,5) & 32 \\
\bottomrule
\end{tabular}
\end{table}

\begin{figure}[H]
    \centering
    \includegraphics[width=0.75\linewidth]{images/master/global_loc_zoom.png}
    \caption{Результат автоматичного пошуку для кадру 2: сітка перевірених фрагментів, фрагменти, що пройшли фільтр (помаранчеві), фрагмент-переможець (синій), контур огляду камери (зелений) і центр кадру (червоний)}
    \label{fig:global-loc}
\end{figure}

В обох випадках переможцем став фрагмент, суміжний з еталонним фрагментом 196, який також містить центр кадру. Оцінка LoFTR відрізняється від калібрувальної оцінки $\approx(1110;\ 2330)$ на $\approx 3$~\px{}, оцінка LightGlue --- на $\approx 19$~\px{}. Усі фрагменти, що пройшли фільтр, перетинаються з місцем зйомки; жоден віддалений фрагмент фільтр не пройшов. Перегляд усієї карти з кешем ознак займає 33~с (85~мс на фрагмент), що прийнятно для початкової ініціалізації, але не для обробки кожного кадру. Це обґрунтовує використання трекера.

\subsection{Локалізація послідовності кадрів і вплив компенсації курсу}
\label{sec:exp-yaw}

Траєкторія польоту являє собою облітання кварталу: на північ, схід, південь і захід. Під час польоту курс змінюється приблизно так: $0^\circ \to 90^\circ \to 175^\circ \to -90^\circ \to -40^\circ$. Щоб оцінити внесок компенсації курсу, трекер (LightGlue + RANSAC) запущено на всіх 95 кадрах з компенсацією та без неї (табл.~\ref{tab:yaw}, рис.~\ref{fig:yaw}).

\begin{table}[H]
\centering
\caption{Вплив компенсації курсу на трекінг послідовності (LightGlue + RANSAC, 95 кадрів)}
\label{tab:yaw}
\begin{tabular}{lcc}
\toprule
& \textbf{Без компенсації} & \textbf{З компенсацією} \\
\midrule
\textsc{initialized} / \textsc{tracking} / \textsc{interpolated} / \textsc{lost} & 1 / 17 / 3 / 74 & 1 / 94 / 0 / 0 \\
Супровід утрачено & на кадрі 19, не відновлено & ні \\
Загальний час обробки & $\approx 41$~хв & 133~с \\
\bottomrule
\end{tabular}
\end{table}

\begin{figure}[H]
    \centering
    \begin{subfigure}{0.49\linewidth}
        \centering
        \includegraphics[width=\linewidth]{images/master/trajectory_no_yaw.png}
        \caption{без компенсації курсу}
    \end{subfigure}
    \hfill
    \begin{subfigure}{0.49\linewidth}
        \centering
        \includegraphics[width=\linewidth]{images/master/trajectory_with_yaw.png}
        \caption{з компенсацією курсу}
    \end{subfigure}
    \caption{Траєкторія, відновлена трекером: зелені точки --- \textsc{tracking}, жовті --- \textsc{interpolated}, червоні хрестики --- кандидати, відкинуті через низький бал або кінематичну перевірку}
    \label{fig:yaw}
\end{figure}

Без компенсації трекер утрачає супровід на першому розвороті (кадр 19, поворот кадру $\approx 50^\circ$). Глобальний пошук із курсом $0^\circ$ після цього вже не знаходить відповідного фрагмента. Кожен кадр у стані \textsc{lost} займає $\approx 32$~с, тому загальний час обробки зростає до 41~хв. З компенсацією всі чотири розвороти відстежено лише за прогнозом курсу; перебір додаткових гіпотез курсу знадобився для одного кадру. Компенсація курсу виявилася необхідною умовою роботи системи на реальному польоті.

\subsection{Порівняння методів LightGlue і LoFTR}
\label{sec:exp-methods}

Усі чотири конфігурації відстежили всю послідовність: 94 кадри зі статусом \textsc{tracking} і один \textsc{initialized}, лише для LightGlue + MAGSAC++ один кадр екстрапольовано. Кінематична перевірка жодного разу не відкинула оцінку, а повторна глобальна ініціалізація не знадобилася. Показники трекінгу та точності наведено в табл.~\ref{tab:tracking} і \ref{tab:accuracy}; відновлені траєкторії разом з еталоном --- на рис.~\ref{fig:traj-gt}, похибки на опорних кадрах --- на рис.~\ref{fig:error}.

\begin{table}[H]
\centering
\caption{Показники трекінгу послідовності з 95 кадрів}
\label{tab:tracking}
\small
\begin{tabular}{lcccccc}
\toprule
\multirow{2}{*}{\textbf{Конфігурація}} & \multicolumn{2}{c}{\textbf{Час кадру, с}} & \textbf{Глобальний} & \multirow{2}{*}{\textbf{Медіана} $N_{\text{in}}$} & \multirow{2}{*}{$\bar{r}$} & \multirow{2}{*}{$\bar{e}$, \px{}} \\
& середній & $P_{90}$ & \textbf{пошук, с} & & & \\
\midrule
LightGlue + RANSAC & 1,08 & 1,40 & 38,3 & 81 & 0,48 & 1,65 \\
LightGlue + MAGSAC++ & 0,78 & 1,02 & 20,1 & 87 & 0,50 & 1,73 \\
LoFTR + RANSAC & 2,72 & 2,93 & 127,9 & 316 & 0,59 & 1,27 \\
LoFTR + MAGSAC++ & 2,27 & 2,49 & 87,5 & 353 & 0,61 & 1,33 \\
\bottomrule
\end{tabular}
\end{table}

\begin{table}[H]
\centering
\caption{Похибка локалізації на 18 надійних опорних кадрах, \px{}}
\label{tab:accuracy}
\small
\begin{tabular}{lccccccc}
\toprule
\textbf{Конфігурація} & \textbf{Середня} & \textbf{Медіана} & \textbf{RMSE} & $P_{90}$ & \textbf{Макс.} & \textbf{S@10} & \textbf{S@25} \\
\midrule
LightGlue + RANSAC & 7,2 & 6,0 & 9,5 & 16,8 & 22,9 & 83\,\% & 100\,\% \\
LightGlue + MAGSAC++ & 7,5 & 5,7 & 10,2 & 16,5 & 27,5 & 78\,\% & 94\,\% \\
LoFTR + RANSAC & 4,9 & 3,1 & 7,3 & 9,0 & 24,2 & 94\,\% & 100\,\% \\
LoFTR + MAGSAC++ & 5,2 & 2,9 & 8,9 & 8,2 & 33,2 & 94\,\% & 94\,\% \\
\bottomrule
\end{tabular}
\end{table}

\begin{figure}[H]
    \centering
    \includegraphics[width=0.8\linewidth]{images/master/trajectories_vs_gt.png}
    \caption{Траєкторії, відновлені чотирма конфігураціями, та еталонні положення опорних кадрів}
    \label{fig:traj-gt}
\end{figure}

\begin{figure}[H]
    \centering
    \includegraphics[width=\linewidth]{images/master/error_over_time.png}
    \caption{Похибка локалізації на опорних кадрах уздовж польоту}
    \label{fig:error}
\end{figure}

\textbf{Точність.}\\
LoFTR забезпечує меншу медіанну похибку, ніж LightGlue: $\approx 3$~\px{} проти $\approx 6$~\px{}. Частка кадрів із похибкою не більше 10~\px{} становить 94\,\% проти 78--83\,\%. У парному порівнянні з однаковим оцінювачем RANSAC LoFTR точніший на 14 з 18 опорних кадрів (знаковий тест, $p = 0{,}031$), з MAGSAC++ --- на 13 з 18 ($p = 0{,}096$). На більшості кадрів різниця між методами не перевищує 1--2~\px{}, тобто є порівнянною з невизначеністю еталона. Перевагу LoFTR формують насамперед кілька складних кадрів:
\begin{itemize}
    \item кадр 1: 19,9~\px{} у LightGlue проти 2,1~\px{} у LoFTR --- LightGlue під час ініціалізації оцінює гомографію площини даху високої будівлі (див. п.~\ref{sec:exp-pair});
    \item кадр 82 поблизу масиву дерев: 15,5~\px{} проти 3,8~\px{};
    \item кадр 90: 7,5~\px{} проти 2,4~\px{}.
\end{itemize}

На кадрі 80 усі чотири конфігурації мають близьку похибку 23--33~\px{}. Збіг похибки незалежних методів вказує на систематичну причину, спільну для всіх. Найімовірніше, це неточність еталона: усі точки розмітки цього кадру розташовані в його лівій третині (дорога й паркінг), права частина кадру зайнята деревами, а центр кадру екстрапольовано. Без кадру 80 середня похибка становить 6,2--6,3~\px{} для LightGlue і 3,6--3,7~\px{} для LoFTR; максимальна похибка LoFTR зменшується до 9,2~\px{}.

\textbf{Кількість відповідностей і швидкодія.}\\
LoFTR знаходить у 3,6--4 рази більше геометрично узгоджених відповідностей (медіана 316--353 проти 81--87) з більшою їх часткою та меншою похибкою репроєкції. Водночас LoFTR у 2,5--2,9 раза повільніший у режимі трекінгу (2,3--2,7~с проти 0,8--1,1~с на кадр) і в 3,3 (RANSAC) та 4,4 (MAGSAC++) раза повільніший під час глобального пошуку, оскільки для нього неможливе кешування ознак фрагментів. Між собою траєкторії конфігурацій відрізняються з медіаною 0,8--1,7~\px{}; найбільші розбіжності (19--43~\px{}) припадають на кадри 79--83 над деревами.

\subsection{Порівняння оцінювачів RANSAC і MAGSAC++}

Розподіл часу обробки кадру в режимі трекінгу наведено на рисунку~\ref{fig:timing}.

\begin{figure}[H]
    \centering
    \includegraphics[width=0.75\linewidth]{images/master/timing_local.png}
    \caption{Час обробки кадру в режимі \textsc{tracking} для чотирьох конфігурацій}
    \label{fig:timing}
\end{figure}

Заміна RANSAC на MAGSAC++ скорочує час кадру на 28\,\% для LightGlue (1,08 $\to$ 0,78~с) і на 17\,\% для LoFTR (2,72 $\to$ 2,27~с), а час глобального пошуку --- у 1,9 і 1,5 раза відповідно. Виграш отримано насамперед на фрагментах, які не містять місця зйомки: на таких фрагментах реалізація RANSAC в OpenCV не може достроково зупинитися і виконує всі 5000 ітерацій, тоді як MAGSAC++ завершується за кілька мілісекунд.

Статистично значущої різниці в точності не виявлено. У парному порівнянні MAGSAC++ точніший за RANSAC на 12 з 18 кадрів для LightGlue ($p = 0{,}24$) і на 9 з 18 для LoFTR ($p = 1{,}0$), медіанна різниця не перевищує 0,2~\px{}. Траєкторії, отримані з різними оцінювачами, відрізняються з медіаною 1,2~\px{} (LightGlue) і 0,8~\px{} (LoFTR). Таким чином, MAGSAC++ дає змогу пришвидшити систему без втрати точності. Конфігурація LightGlue + MAGSAC++ --- єдина, що забезпечує обробку кадру не довше 1~с (середнє 0,78~с, $P_{90}$ 1,02~с).

\subsection{Стійкість до спотворень зображення}
\label{sec:exp-robust}

Для 19 кадрів (кожен п'ятий кадр послідовності) штучно застосовано спотворення до оригінального кадру $1920 \times 1080$:
\begin{itemize}
    \item зниження яскравості $I' = \alpha I$;
    \item зниження контрасту $I' = \alpha (I - \mu) + \mu$, де $\mu$ --- середнє значення каналу;
    \item адитивний гауссів шум;
    \item лінійне розмиття руху;
    \item поворот кадру без компенсації;
    \item помилка у заданому масштабі $s_f$.
\end{itemize}
Пул фрагментів і компенсація курсу були такими самими, як у трекері на відповідному кадрі. Для кожного кадру обчислено відношення $N_{\text{in}}$ до значення для неспотвореного кадру та зміщення знайденого центру. Результати наведено в табл.~\ref{tab:robust} і на рис.~\ref{fig:robust}.

\begin{table}[H]
\centering
\caption{Стійкість до спотворень (19 кадрів): частка локалізованих кадрів і медіанне відношення $N_{\text{in}} / N_{\text{in}}^{\text{без спотв.}}$}
\label{tab:robust}
\small
\begin{tabular}{lcccc}
\toprule
\multirow{2}{*}{\textbf{Спотворення}} & \multicolumn{2}{c}{\textbf{LightGlue}} & \multicolumn{2}{c}{\textbf{LoFTR}} \\
\cmidrule(lr){2-3}\cmidrule(lr){4-5}
& локалізовано & відношення & локалізовано & відношення \\
\midrule
яскравість $\alpha = 0{,}7$ / $0{,}5$ & 100 / 100\,\% & 1,00 / 1,01 & 100 / 100\,\% & 0,97 / 0,97 \\
контраст $\alpha = 0{,}7$ / $0{,}5$ & 100 / 100\,\% & 1,02 / 0,99 & 100 / 100\,\% & 1,00 / 0,97 \\
шум $\sigma = 10$ / $25$ & 95 / 95\,\% & 0,99 / 0,92 & 100 / 95\,\% & 1,02 / 0,95 \\
розмиття 5 / 9 пікселів & 100 / 100\,\% & 1,02 / 1,02 & 100 / 100\,\% & 1,02 / 1,04 \\
поворот $15^\circ$ / $30^\circ$ & 95 / 95\,\% & 1,03 / 0,99 & 95 / 95\,\% & 0,82 / 0,52 \\
поворот $45^\circ$ & 53\,\% & 0,35 & 11\,\% & 0 \\
поворот $90^\circ$ & 0\,\% & 0 & 0\,\% & 0 \\
масштаб $\times 0{,}8$ / $\times 1{,}25$ & 100 / 95\,\% & 0,68 / 1,29 & 100 / 100\,\% & 0,71 / 1,05 \\
\bottomrule
\end{tabular}
\end{table}

\begin{figure}[H]
    \centering
    \includegraphics[width=\linewidth]{images/master/robustness.png}
    \caption{Стійкість методів до спотворень кадру}
    \label{fig:robust}
\end{figure}

Зміни яскравості й контрасту, гауссів шум і розмиття в досліджених межах практично не впливають на результат: кількість геометрично узгоджених відповідностей зберігається, а медіанне зміщення центру не перевищує 1,5~\px{}. Слабкий вплив розмиття частково пояснюється тим, що після масштабування кадру з коефіцієнтом 0,27 ядро 9 пікселів зменшується до $\approx 2{,}4$ пікселя.

Визначальним чинником виявився поворот кадру. LightGlue зберігає працездатність до $30^\circ$, а LoFTR втрачає половину відповідностей уже при $30^\circ$. При $45^\circ$ локалізовано лише 53\,\% (LightGlue) і 11\,\% (LoFTR) кадрів, при $90^\circ$ --- жодного. Ці результати кількісно обґрунтовують компенсацію курсу, описану в п.~\ref{sec:tracking}. Похибка масштабу на 20--25\,\% не призводить до втрати локалізації, але зменшення заданого масштабу на 20\,\% скорочує кількість відповідностей приблизно на 30\,\%.

\subsection{Аналіз проблемних ділянок місцевості}
\label{sec:exp-terrain}

Для кожного кадру за контуром огляду обчислено дві ознаки ділянки ортофотоплану:
\begin{itemize}
    \item частку рослинності за індексом Excess Green $\text{ExG} = 2g - r - b$ у хроматичних координатах \cite{exg} (піксель вважається рослинністю, якщо $\text{ExG} > 0{,}05$);
    \item текстурованість --- середній модуль градієнта Собеля яскравості.
\end{itemize}
Залежність кількості геометрично узгоджених відповідностей від номера кадру та частки рослинності наведено на рисунку~\ref{fig:terrain}, кореляції --- у табл.~\ref{tab:terrain}.

\begin{figure}[H]
    \centering
    \includegraphics[width=\linewidth]{images/master/inliers_vs_terrain.png}
    \caption{Кількість геометрично узгоджених відповідностей уздовж польоту (вгорі) і частка рослинності під кадром (внизу); сірим виділено кадри 81--84}
    \label{fig:terrain}
\end{figure}

\begin{table}[H]
\centering
\caption{Зв'язок кількості геометрично узгоджених відповідностей з ознаками місцевості (кореляція Спірмена $\rho$, RANSAC)}
\label{tab:terrain}
\begin{tabular}{lcccc}
\toprule
\textbf{Метод} & $\rho$ (рослинність) & $\rho$ (текстура) & \textbf{Медіана} $N_{\text{in}}$, кадри 81--84 & \textbf{Медіана} $N_{\text{in}}$, усі \\
\midrule
LightGlue & $+0{,}37$ & $+0{,}65$ & 27 & 81 \\
LoFTR & $-0{,}29$ & $+0{,}29$ & 138 & 316 \\
\bottomrule
\end{tabular}
\end{table}

На ділянці кадрів 81--84 (паркінг і канал, оточені деревами) кількість геометрично узгоджених відповідностей падає втричі для LightGlue і вдвічі для LoFTR. Саме тут методи найбільше розходяться між собою. Водночас частка рослинності під кадром сама по собі не пояснює це падіння: для LightGlue кореляція з нею навіть додатна, оскільки дачні ділянки з будівлями мають високу частку рослинності, але багаті на деталі. Кращим предиктором для LightGlue є текстурованість ($\rho = 0{,}65$); для LoFTR зв'язок з обома ознаками слабкий. Імовірною причиною падіння на кадрах 81--84 є сезонна різниця крон листяних дерев (на відео --- без листя, на карті --- з листям), однак підтвердження цієї гіпотези потребує окремого дослідження.

\subsection{Дослідження вдосконалень системи}
\label{sec:exp-improve}

\textbf{Методика.}\\
Аналіз результатів основної конфігурації виявив кілька слабких місць: кадр не вміщується у фрагмент 512~\px{}, вибір фрагмента залежить від того, як кадр ліг на сітку, гомографія може описувати площину даху, а трекінг LoFTR триває 2,3--2,7~с на кадр. Для їх усунення запропоновано вдосконалення, описані в розділі~\ref{sec:implementation}:
\begin{itemize}
    \item альтернативне розбиття карти: фрагменти $520 \times 520$ з перекриттям 33\,\% (216 фрагментів);
    \item штраф за відхилення масштабу в балі ранжування, $\sigma_s = 0{,}15$ (формула~\ref{eq:scale-prior});
    \item вікно карти за прогнозом у стані \textsc{tracking}, запас $m_w = 128$~\px{} (рис.~\ref{fig:window});
    \item уточнення положення після глобального пошуку, запас 64~\px{}, допустиме зміщення 30~\px{};
    \item згладжування траєкторії фільтром Калмана та згладжувачем RTS (п.~\ref{sec:kalman}).
\end{itemize}
Спершу оцінено вплив розбиття: основна конфігурація без вдосконалень (\emph{базовий варіант}) запущена на фрагментах $520 \times 520$. Далі на тих самих фрагментах для кожної з чотирьох конфігурацій запущено ще чотири варіанти: з кожним із трьох вдосконалень трекера окремо та з усіма разом --- разом 20 прогонів трекінгу на 95 кадрах. Згладжування не є окремим прогоном: його застосовано до траєкторії кожного прогону. Точність оцінено за тими самими 18 надійними опорними кадрами, що й у п.~\ref{sec:exp-methods}; основні результати й дані не змінювалися.

\textbf{Рівень шуму порівнянь.}\\
Перш ніж порівнювати варіанти, потрібно оцінити, яку різницю слід вважати значущою. Уточнення змінює лише перший кадр, однак далі трекер прогнозує дещо інші положення, обирає інші фрагменти й отримує інші оцінки. На решті 17 опорних кадрів траєкторії з уточненням і без нього відрізняються з медіаною 0,2--0,9~\px{} (окремі кадри --- до 20~\px{}). Отже, зміна медіанної похибки на $\approx 1$~\px{} на вибірці з 18 кадрів може бути спричинена самою чутливістю трекера до незначних змін. Тому, як і раніше, для кожного порівняння наведено парний знаковий тест.

\textbf{Вплив розбиття карти.}\\
Результати базового варіанта на двох розбиттях наведено в табл.~\ref{tab:tiling}.

\begin{table}[H]
\centering
\caption{Основна конфігурація на двох розбиттях карти (95 кадрів, 18 надійних опорних кадрів)}
\label{tab:tiling}
\small
\begin{tabular}{lcccccc}
\toprule
\multirow{2}{*}{\textbf{Конфігурація}} & \multicolumn{2}{c}{\textbf{Глобальний пошук, с}} & \multicolumn{2}{c}{\textbf{Медіана похибки, \px{}}} & \multicolumn{2}{c}{\textbf{S@10}} \\
\cmidrule(lr){2-3}\cmidrule(lr){4-5}\cmidrule(lr){6-7}
& 512, 50\,\% & 520, 33\,\% & 512, 50\,\% & 520, 33\,\% & 512, 50\,\% & 520, 33\,\% \\
\midrule
LightGlue + RANSAC   & 38,3  & 19,6 & 6,0 & 4,6 & 83\,\% & 89\,\% \\
LightGlue + MAGSAC++ & 20,1  & 11,1 & 5,7 & 3,2 & 78\,\% & 94\,\% \\
LoFTR + RANSAC       & 127,9 & 73,7 & 3,1 & 3,9 & 94\,\% & 94\,\% \\
LoFTR + MAGSAC++     & 87,5  & 52,9 & 2,9 & 3,4 & 94\,\% & 94\,\% \\
\bottomrule
\end{tabular}
\end{table}

Завдяки меншій кількості фрагментів (216 замість 384) глобальний пошук пришвидшився в 1,7--2,0 раза. Час кадру в режимі трекінгу для LightGlue майже не змінився (0,72--1,06~с), а для LoFTR зріс на 5--14\,\% (2,38--3,09~с) через більший розмір фрагментів. Точність LightGlue зросла: медіанна похибка зменшилася з 5,7--6,0 до 3,2--4,6~\px{}, $P_{90}$ --- з 16,5--16,8 до 8,9--11,6~\px{}. Однак у парному порівнянні нове розбиття краще лише на 11--12 з 18 кадрів ($p = 0{,}24$--$0{,}48$): виграш зосереджений на кількох складних кадрах. Зокрема, на кадрі 1 LightGlue більше не оцінює площину даху (похибка 0,7--2,5~\px{} замість 19--20~\px{}), а на кадрі 95 з MAGSAC++ похибка зменшилася з 12,1 до 3,1~\px{}. Водночас на кадрі 82 над деревами з RANSAC вона зросла з 15,5 до 25,4~\px{}. Для LoFTR зміна точності незначуща: нове розбиття краще на 6--8 з 18 кадрів, медіанна парна різниця становить $+0{,}1$~\px{}.

\textbf{Порівняння варіантів системи.}\\
Результати варіантів на фрагментах $520 \times 520$ наведено в табл.~\ref{tab:variants} і на рис.~\ref{fig:variants}. В усіх 20 прогонах трекер відстежив усю послідовність без утрати супроводу; лише в одному прогоні (LightGlue + MAGSAC++, усі вдосконалення) один кадр екстрапольовано.

\begin{table}[H]
\centering
\caption{Варіанти системи на фрагментах $520 \times 520$: медіанна похибка на опорних кадрах, \px{} / середній час кадру в режимі \textsc{tracking}, с}
\label{tab:variants}
\small
\begin{tabular}{lcccc}
\toprule
\textbf{Варіант} & \textbf{\shortstack{LightGlue\\+ RANSAC}} & \textbf{\shortstack{LightGlue\\+ MAGSAC++}} & \textbf{\shortstack{LoFTR\\+ RANSAC}} & \textbf{\shortstack{LoFTR\\+ MAGSAC++}} \\
\midrule
базовий               & 4,58 / 1,06 & 3,24 / 0,72 & 3,88 / 3,09 & 3,44 / 2,38 \\
штраф за масштаб      & 5,09 / 1,04 & 4,47 / 0,74 & 3,88 / 3,16 & 3,58 / 2,40 \\
вікно за прогнозом    & 3,50 / 0,31 & 3,80 / 0,25 & 3,23 / 0,46 & 3,24 / 0,47 \\
уточнення             & 5,03 / 1,01 & 5,16 / 0,70 & 3,43 / 3,06 & 4,05 / 2,43 \\
усі вдосконалення     & 5,01 / 0,27 & 4,35 / 0,27 & 3,18 / 0,49 & 2,71 / 0,46 \\
\bottomrule
\end{tabular}
\end{table}

\begin{figure}[H]
    \centering
    \includegraphics[width=0.9\linewidth]{images/master/improvements_variants.png}
    \caption{Медіанна похибка на опорних кадрах (вгорі) і середній час кадру в режимі \textsc{tracking} (внизу) для основної конфігурації на двох розбиттях карти та варіантів системи на фрагментах $520 \times 520$; штрихуванням позначено MAGSAC++}
    \label{fig:variants}
\end{figure}

\textbf{Вікно карти за прогнозом.}\\
Вікно за прогнозом виявилося найефективнішим вдосконаленням. З вікна прийнято 88--90 із 94 кадрів трекінгу для LightGlue і 93--94 для LoFTR. Резервний пошук у фрагментах знадобився на різких розворотах (зміна курсу до $42^\circ$ між кадрами) і над деревами. Час кадру скоротився в 2,9--3,4 раза для LightGlue (до 0,25--0,31~с) і в 5,1--6,7 раза для LoFTR (до 0,46--0,47~с; $P_{90}$ --- 0,50--0,52~с). Обробка всієї послідовності LoFTR + MAGSAC++ скоротилася з 277 до 97~с. Оскільки кадр потрапляє у вікно повністю, медіанна кількість геометрично узгоджених відповідностей зросла з 70--74 до 84--85 для LightGlue і з 258--294 до 307--323 для LoFTR. Точність статистично значуще не змінилася: у парному порівнянні з базовим варіантом $p \ge 0{,}14$, медіанна різниця не перевищує 0,25~\px{}. Таким чином, вікно за прогнозом пришвидшує трекінг без втрати точності, а LoFTR із вікном обробляє кадр швидше, ніж LightGlue у базовому варіанті. Обмеження в 1~с на кадр, яке в основній конфігурації виконувала лише LightGlue + MAGSAC++, тепер виконують усі чотири конфігурації.

\textbf{Штраф за відхилення масштабу.}\\
Штраф не дав вимірного ефекту. Для LoFTR + RANSAC траєкторія не змінилася взагалі, для LoFTR + MAGSAC++ змінилися оцінки лише на окремих кадрах, а для LightGlue медіанна похибка зросла на 0,5--1,2~\px{} за статистично незначущої парної різниці ($p = 0{,}30$--$0{,}61$). Ситуація, яка мотивувала штраф (вибір фрагмента з гомографією площини даху), на розбитті $520 \times 520$ у базовому варіанті не виникла: на кадрі 1 похибка LightGlue і без штрафу становить 0,7--2,5~\px{}. Штраф доцільно розглядати як запобіжник, ефект якого потребує перевірки на польотах із більшою кількістю високих будівель.

\textbf{Уточнення положення після глобального пошуку.}\\
Уточнення спрацювало в усіх прогонах (на першому кадрі), але дало неоднозначний результат. Для LoFTR похибка кадру 1 зменшилася (з 2,1 до 0,9~\px{} і з 1,5 до 1,0~\px{}), а для LightGlue --- зросла з 0,7 до 20,4~\px{} (RANSAC) і з 2,5 до 10,2~\px{} (MAGSAC++). Причиною є площина даху: ділянка карти навколо контуру кадру містить дах високої будівлі повністю, і після уточнення масштаб гомографії LightGlue зменшився з 0,275 до 0,206 (RANSAC) і з 0,269 до 0,248 (MAGSAC++), тоді як для LoFTR він залишився рівним 0,27. Кількість узгоджених відповідностей при цьому зросла (з 94 до 103 і з 98 до 107), тому умови прийняття уточнення (не менше відповідностей, зміщення до 30~\px{}) від такої помилки не захищають: зміщення центру в 10--20~\px{} не перевищує допуску. Через зміну першого кадру змінюється і решта траєкторії, тому медіанна похибка LightGlue з уточненням зросла на 0,5--1,9~\px{}. Отже, у запропонованому вигляді уточнення для LightGlue недоцільне. Його можна вдосконалити, поєднавши зі штрафом за масштаб у критерії прийняття.

\textbf{Усі вдосконалення разом.}\\
Найменшу медіанну похибку в серії отримано для LoFTR + MAGSAC++ з усіма вдосконаленнями: 2,71~\px{} ($P_{90}$ 8,5~\px{}, 94\,\% кадрів із похибкою до 10~\px{}) за часу 0,46~с на кадр. Втім, порівняно з базовим варіантом на тих самих фрагментах ця різниця статистично незначуща (краще на 11 з 18 кадрів, $p = 0{,}48$). Для LightGlue поєднання всіх вдосконалень дало гірший результат, ніж лише вікно (4,35--5,01~\px{} проти 3,50--3,80~\px{}), через описаний вплив уточнення на перший кадр.

\textbf{Згладжування траєкторії.}\\
Результати згладжування траєкторій базового варіанта наведено в табл.~\ref{tab:smoothing}.

\begin{table}[H]
\centering
\caption{Вплив згладжування на похибку на опорних кадрах (базовий варіант, фрагменти $520 \times 520$): медіана / RMSE, \px{}; у дужках --- кількість кадрів, на яких похибка менша, ніж без згладжування}
\label{tab:smoothing}
\small
\begin{tabular}{lccc}
\toprule
\textbf{Конфігурація} & \textbf{Без згладжування} & \textbf{Фільтр Калмана} & \textbf{Згладжувач RTS} \\
\midrule
LightGlue + RANSAC   & 4,58 / 8,76 & 4,14 / 8,58 (9/18)  & 3,97 / 8,12 (9/18) \\
LightGlue + MAGSAC++ & 3,24 / 6,26 & 3,16 / 6,21 (12/18) & 3,28 / 6,30 (8/18) \\
LoFTR + RANSAC       & 3,88 / 8,30 & 3,55 / 8,34 (5/18)  & 3,63 / 8,93 (9/18) \\
LoFTR + MAGSAC++     & 3,44 / 9,88 & 3,44 / 9,74 (9/18)  & 3,34 / 9,13 (11/18) \\
\bottomrule
\end{tabular}
\end{table}

Згладжування змінює медіанну похибку не більше ніж на 0,6~\px{} і в різних напрямках, а кількість кадрів, на яких похибка зменшилася, коливається навколо половини (5--12 з 18). Отже, на опорних кадрах похибка визначається не випадковим <<тремтінням>> траєкторії, яке усуває фільтр, а систематичними чинниками: неплоскою сценою, ділянками над деревами та невизначеністю самого еталона (зокрема на кадрі 80, де похибка всіх варіантів становить 18--52~\px{}). Згладжування доцільне для отримання плавної траєкторії, однак точність локалізації в досліджених умовах воно не підвищує.

\subsection{Перевірка системи на іншій місцевості}
\label{sec:exp-flight2}

\textbf{Мета.}\\
Перший політ виконано над міською та дачною забудовою з великою кількістю чітких деталей. Щоб з'ясувати, чи переносяться отримані результати на інші умови, систему перевірено на другому польоті над сільськогосподарською місцевістю. Алгоритм при цьому не змінювався. На відміну від першого польоту, для другого є журнал бортової навігації. Це дає незалежний еталон для кожного кадру і дає змогу виразити похибку в метрах (п.~\ref{sec:datasets}).

\textbf{Вхідні дані.}\\
Відеозапис зроблено квадрокоптером DJI Mavic Mini 24.11.2025 над полями поблизу автомобільної траси:
\begin{itemize}
    \item роздільність $960 \times 540$, частота 30~кадр/с, тривалість 335~с;
    \item камера спрямована вертикально вниз (кут нахилу підвісу $-90^\circ$);
    \item висота за барометром 99,7--101,8~м відносно точки зльоту;
    \item довжина маршруту $\approx 2{,}2$~км, швидкість $\approx 6{,}3$~м/с.
\end{itemize}
З відео вилучено 335 кадрів із частотою 1~кадр/с (рис.~\ref{fig:f2-track}).

Ортофотопланом слугує супутниковий знімок сервісу Esri World Imagery \cite{esri}, експортований у середовищі QGIS \cite{qgis} у проєкції UTM (зона 34N) разом із файлом прив'язки. Знімок має розмір $4730 \times 4029$ пікселів і розмір пікселя 0,30~м, тобто покриває ділянку $\approx 1{,}4 \times 1{,}2$~км. Його зроблено влітку: поля на ньому зелені або з посівами, а на відео --- листопадова рілля й дерева без листя (рис.~\ref{fig:f2-frames}). Перевірено також осінній архівний знімок Google Earth, суміщений із картою за відповідностями (залишкова похибка суміщення $\approx 1{,}6$~м). Однак через нижчу роздільність він дав менше відповідностей: на кадрі 100 LoFTR знайшов 22 геометрично узгоджені відповідності проти 91 на літньому знімку. Тому далі використано літній знімок.

\begin{figure}[H]
    \centering
    \includegraphics[width=0.75\linewidth]{images/master/flight2_gt_track.png}
    \caption{Траєкторія другого польоту за даними GNSS на супутниковому знімку (номери кадрів через кожні 30~с); траєкторія використовується лише як еталон}
    \label{fig:f2-track}
\end{figure}

\begin{figure}[H]
    \centering
    \includegraphics[width=\linewidth]{images/master/flight2_frames_vs_map.png}
    \caption{Кадри другого польоту (вгорі) і ділянки карти $180 \times 180$~м навколо еталонного положення (внизу, хрестик); кадри не повернуті до орієнтації карти}
    \label{fig:f2-frames}
\end{figure}

\textbf{Параметри системи.}\\
Параметри задано апріорно, без використання еталона.
\begin{itemize}
    \item \textbf{Масштаб кадру} $s_f = 0{,}52$. Ширина кадру на місцевості за висоти $\approx 101$~м і горизонтального кута огляду камери $\approx 73^\circ$ становить $\approx 150$~м, тобто $\approx 499$~\px{} карти, звідки $s_f = 499 / 960 \approx 0{,}52$. Оцінка масштабу лише за зображеннями (глобальний пошук із перебором курсів на кадрі 320) дала близькі значення 0,44 (LightGlue) і 0,57 (LoFTR), тому апріорне значення збережено.
    \item \textbf{Фрагменти} $496 \times 496$ з перекриттям 33\,\%, разом 168 фрагментів. Сторона фрагмента дорівнює ширині контуру огляду ($\approx 499 \times 281$~\px{}), округленій до кратного 8.
    \item \textbf{Решта параметрів} такі самі, як у п.~\ref{sec:exp-improve}. Використано оцінювач MAGSAC++ і вікно карти за прогнозом.
\end{itemize}

\textbf{Результати без компаса.}\\
Трекінг усього польоту з LightGlue + MAGSAC++ (основна конфігурація) виявився значно гіршим, ніж на першому польоті:
\begin{itemize}
    \item 309 із 335 кадрів залишилися без положення (\textsc{lost});
    \item першу прив'язку отримано лише на кадрі 303 біля траси;
    \item положення визначено на кадрах 303--323 і 331--335: 23 оцінки з медіанною похибкою 5,2~м і максимальною 9,1~м.
\end{itemize}
Отже, хибних прив'язок система не дала: там, де вона видала положення, воно правильне, але таких кадрів лише 7\,\%. Глобальний пошук триває $\approx 5{,}8$~с на кадр, а обробка всього польоту зайняла 30,5~хв.

LoFTR не кешує ознак карти, тому його глобальний пошук значно повільніший: 50--70~с на кадр. Через це LoFTR запускався лише на відрізку кадрів 85--135, де маршрут проходить над дорогою й лісосмугою. На цьому відрізку без компаса не локалізовано жодного з 51 кадру.

\textbf{Аналіз причин.}\\
Щоб відокремити помилки пошуку від меж самого зіставлення, виконано діагностику з підказкою. Кадр, повернутий на еталонний курс, зіставлявся з ділянкою карти навколо еталонного положення. Цей режим використовує еталон і застосовувався лише для аналізу, а не для налаштування системи. Результати діагностики:
\begin{itemize}
    \item LoFTR надійно зіставляє кадри над дорогою й лісосмугою (кадри 30 і 100: до 91 узгодженої відповідності, похибка 1--2~м) і біля траси (кадр 320); кадр 60 --- слабко;
    \item LightGlue зіставляє лише кадри біля траси (300, 320);
    \item над серединою полів (кадри 140--260) обидва методи знаходять лише 0--7 узгоджених відповідностей навіть за відомих положення й курсу;
    \item масштаб гомографії змінювався в межах 0,47--0,62.
\end{itemize}

Таким чином, низьку частку локалізованих кадрів зумовлюють такі причини:
\begin{enumerate}
    \item \textbf{Сезонна різниця.} Однотонні зелені поля на карті на кадрах перетворилися на ріллю з борозенками, а крони дерев --- на голі гілки (рис.~\ref{fig:f2-frames}).
    \item \textbf{Однорідна місцевість.} Кадр охоплює $\approx 150 \times 85$~м, і над серединою поля в ньому немає жодного стабільного орієнтира. Спільних ознак кадру й карти тут немає, і цю межу не можна усунути налаштуванням параметрів.
    \item \textbf{Рельєф.} Висота за барометром відраховується від точки зльоту, а рельєф під апаратом змінюється. Тому фактичний масштаб кадру відрізняється від прийнятого до $\approx 20$\,\%. Як показано в п.~\ref{sec:exp-robust}, така різниця ще допустима, але зменшує кількість відповідностей.
    \item \textbf{Невідомий курс під час глобального пошуку.} На відрізку 85--135 кадр повернутий відносно карти на $\approx 25^\circ$. Найближча гіпотеза циклу $0^\circ, 90^\circ, 180^\circ, -90^\circ$ відрізняється від нього на $\approx 25^\circ$. За слабкої текстури цього достатньо, щоб відповідності не пройшли фільтр: уже при $30^\circ$ LoFTR втрачає половину відповідностей (п.~\ref{sec:exp-robust}).
\end{enumerate}
Перші три причини є властивостями даних. Четверту можна усунути додатковим датчиком, тобто компасом.

\textbf{Експеримент із компасом.}\\
У трекер додано курс із компаса (п.~\ref{sec:tracking}, формула~\ref{eq:compass}). Окремого магнітометра в журналі немає, тому показ компаса імітовано курсом камери із журналу польоту DJI. Цей курс автопілот обчислює з магнітометра та інерціальних датчиків. Координати з журналу системі не передаються.

Поправка $\Delta = 2{,}28^\circ$ дорівнює куту зближення меридіанів сітки UTM 34N у районі польоту. Курс у журналі вже відраховано від географічної півночі, тому магнітне схилення не додається. Поправку обчислено з координат карти, а не підібрано за результатами. Залишкова розбіжність курсу компаса з кутом, оціненим системою з гомографії, становить $\approx 3^\circ$.

Курс у журналі може уточнюватися автопілотом із використанням GNSS, тож у реальних умовах без GNSS компас буде менш точним. Тому система перевірялася у двох варіантах:
\begin{itemize}
    \item з показом компаса, виправленим лише на $\Delta$;
    \item з навмисно погіршеним компасом: стале зміщення $+10^\circ$ і гауссів шум із середньоквадратичним відхиленням $5^\circ$.
\end{itemize}
Результати наведено в табл.~\ref{tab:flight2} і на рис.~\ref{fig:f2-loftr}, \ref{fig:f2-lg}.

\begin{table}[H]
\centering
\caption{Другий політ: локалізація без компаса і з компасом (MAGSAC++; похибка --- на кадрах зі статусами \textsc{tracking} та \textsc{initialized} відносно еталона GNSS)}
\label{tab:flight2}
\small
\begin{tabular}{p{0.25\linewidth}cccccc}
\toprule
\textbf{Варіант} & \textbf{\shortstack{Кадрів\\усього}} & \textbf{\shortstack{Знайдено\\(+ екстрап.)}} & \textbf{\shortstack{Медіана,\\м}} & \textbf{\shortstack{$P_{90}$,\\м}} & \textbf{\shortstack{Макс.,\\м}} & \textbf{\shortstack{Час,\\с/кадр}} \\
\midrule
\multicolumn{7}{l}{\textit{LoFTR, відрізок кадрів 85--135, вікно за прогнозом}} \\
без компаса                 & 51 & 0 (0)   & ---  & ---  & ---  & 41,2 \\
компас                      & 51 & 30 (6)  & 2,2  & 4,6  & 7,2  & 29,9 \\
компас $+10^\circ$, шум $5^\circ$ & 51 & 27 (8)  & 2,2  & 3,4  & 8,3  & 21,9 \\
\midrule
\multicolumn{7}{l}{\textit{LightGlue, увесь політ}} \\
без компаса, без вікна      & 335 & 23 (3) & 5,2 & 6,8 & 9,1  & 5,5 \\
компас $+10^\circ$, шум $5^\circ$, вікно & 335 & 43 (5) & 5,2 & 6,1 & 25,9 & 3,8 \\
\bottomrule
\end{tabular}
\end{table}

\begin{figure}[H]
    \centering
    \includegraphics[width=0.7\linewidth]{images/master/flight2_loftr_compass.png}
    \caption{Траєкторія LoFTR + MAGSAC++ з компасом на відрізку кадрів 85--135: зеленим --- кадри \textsc{tracking}, синім --- ініціалізація, жовтим --- екстрапольовані кадри, червоними хрестиками --- відкинуті кінематичною перевіркою оцінки; білим --- контури огляду окремих кадрів}
    \label{fig:f2-loftr}
\end{figure}

\begin{figure}[H]
    \centering
    \includegraphics[width=0.85\linewidth]{images/master/flight2_lightglue_compass.png}
    \caption{Траєкторія LightGlue + MAGSAC++ з погіршеним компасом ($+10^\circ$, шум $5^\circ$) на всьому польоті: положення визначено на ділянках кадрів 72--85 і 302--335 (пряма лінія між ними з'єднує сусідні оцінки і не є оцінкою траєкторії)}
    \label{fig:f2-lg}
\end{figure}

\textbf{LoFTR з компасом.}\\
Найбільший ефект компас дав для LoFTR. На відрізку 85--135, де без компаса не локалізовано жодного кадру, з компасом визначено положення 30 кадрів (59\,\%) і ще 6 екстрапольовано. Медіанна похибка становить 2,2~м, максимальна --- 7,2~м. Перша прив'язка відбулася на кадрі 92 з похибкою 1,3~м. Далі трекер супроводжував кадри 92--115, на кадрах 116--122 втратив супровід, а з кадру 123 знову прив'язався і супроводжував кадри до кінця відрізка.

Погіршений компас дав практично такий самий результат: 27 кадрів, медіана 2,2~м, максимум 8,3~м. Похибка компаса $10$--$15^\circ$ вкладається в допуск зіставлення до повороту, а точний кут визначається гомографією. Тому до точності компаса система невимоглива.

Середній час кадру (22--30~с) визначається глобальним пошуком у стані \textsc{lost} (50--68~с на кадр без кешу ознак), а не трекінгом.

\textbf{LightGlue з компасом.}\\
На всьому польоті з погіршеним компасом LightGlue визначив положення 43 кадрів замість 23. Компас дав такі зміни:
\begin{itemize}
    \item знайдено нову ділянку кадрів 72--85 на межі полів (перша прив'язка на кадрі 72 з похибкою 1,3~м);
    \item ділянку біля траси підхоплено на кадрі 302 і супроводжено без розриву до кінця польоту.
\end{itemize}
Одна оцінка (кадр 322) має похибку 25,9~м, решта --- не більше 10~м.

Прогін без компаса виконано без вікна за прогнозом, тому зменшення часу кадру частково зумовлене вікном. Однак нові ділянки знайдено глобальним пошуком, на який вікно не впливає. Отже, приріст кількості знайдених кадрів дав саме компас. Над серединою полів (кадри 86--301) LightGlue не локалізує жодного кадру і з компасом, що узгоджується з результатами діагностики.

\textbf{Обмеження експерименту.}
\begin{itemize}
    \item Відрізок 85--135 для LoFTR обрано за наявністю орієнтирів, тому частка локалізованих кадрів на ньому не характеризує політ загалом. На всьому польоті вона була б нижчою, оскільки середина полів не локалізується жодним методом.
    \item Показ компаса імітовано курсом автопілота, а не сирим магнітометром. Стійкість до неточного компаса перевірено зміщенням і шумом, але не магнітними завадами, які в реальному польоті можуть давати й більші відхилення.
    \item Еталон GNSS має власну похибку $\approx 2$--$3$~м, тож похибки локалізації на рівні 2~м близькі до точності самого еталона.
\end{itemize}


\subsection{Обмеження розробленої системи}

Під час експериментів виявлено такі обмеження:
\begin{itemize}
    \item \textbf{Плоска модель сцени.} Гомографія описує лише площину. Високі будівлі спричиняють паралакс, і RANSAC може оцінити площину даху замість площини землі (кадр 1 у LightGlue, кадр 2 у LoFTR). Уточнення положення на ділянці карти, що повністю містить дах, цей ефект для LightGlue підсилює, а штраф за масштаб на досліджуваному польоті його не усунув.
    \item \textbf{Розмір фрагмента.} Контур огляду кадру ($\approx 518 \times 292$~\px{}) ширший за фрагмент 512~\px{}, тому жоден фрагмент не містить кадр повністю. У режимі трекінгу це обмеження знімає вікно за прогнозом, однак глобальний пошук і резервний пошук на розворотах і далі виконуються на фрагментах сітки.
    \item \textbf{Фіксований масштаб.} Масштаб кадру $s_f$ відкалібровано на одному кадрі, і він вважається сталим. Під час польоту оцінений масштаб гомографії змінювався в межах 0,20--0,32 (медіана 0,269), проте в трекері він не уточнюється. На другому польоті масштаб змінювався в межах 0,47--0,62 через рельєф, оскільки висота за барометром відраховується від точки зльоту.
    \item \textbf{Чутливість до повороту.} Без компенсації курсу обидва методи непрацездатні при повороті кадру на $45^\circ$ і більше. Без компаса глобальний пошук перебирає курси з кроком $90^\circ$, і на слабкотекстурованій місцевості розбіжності $\approx 25^\circ$ вже достатньо для втрати локалізації. Курс із компаса цю проблему усуває.
    \item \textbf{Сезонна різниця та однорідна місцевість.} На другому польоті (літній знімок, листопадове відео) над серединою полів кадр і карта не мають спільних ознак: обидва методи знаходять 0--7 узгоджених відповідностей навіть за відомих положення й курсу. Система локалізується лише за орієнтирами (дороги, лісосмуги, траса) і на однорідних полях положення не визначає.
    \item \textbf{Швидкодія.} Повний перебір карти займає 20--128~с залежно від конфігурації (11--74~с на розбитті $520 \times 520$). В основній конфігурації обмеження 1~с на кадр у режимі трекінгу на використаному GPU виконує лише LightGlue + MAGSAC++; з вікном за прогнозом його виконують усі конфігурації (0,25--0,47~с), але на різких розворотах кадр обробляється резервним пошуком і займає стільки ж часу, скільки в основній конфігурації.
    \item \textbf{Чутливість порівнянь.} Трекер чутливий до незначних змін: зміна оцінки лише першого кадру змінює подальшу траєкторію на 0,2--0,9~\px{} у медіані. Тому на вибірці з 18 опорних кадрів відмінності між варіантами системи на рівні $\approx 1$~\px{} не є статистично значущими.
    \item \textbf{Одиниці вимірювання.} Ортофотоплан першого польоту не має просторової прив'язки, тому похибки для нього визначено в пікселях карти. Для другого польоту похибки наведено в метрах, але еталон GNSS має власну похибку $\approx 2$--$3$~м.
    \item \textbf{Обсяг даних і еталона.} Основне дослідження виконано на одному польоті (95 кадрів) і 18 опорних кадрах. Опорні кадри обрано за наявністю наземних орієнтирів, тому міські ділянки в еталоні представлені ширше, ніж у польоті загалом. Другий політ (335 кадрів) використано для перевірки на іншій місцевості; LoFTR на ньому через тривалість глобального пошуку досліджено лише на відрізку з 51 кадру.
    \item \textbf{Курс із компаса.} Компас імітовано курсом автопілота з журналу польоту, який може уточнюватися за допомогою GNSS. Стійкість до похибки компаса перевірено зміщенням $10^\circ$ і шумом $5^\circ$, але не реальними магнітними завадами.
\end{itemize}

\subsection{Висновки за результатами експериментів}

\begin{enumerate}
    \item Розроблена система автоматично знаходить відповідний фрагмент ортофотоплану для кадру БпЛА без апріорної інформації про місце зйомки та відновлює траєкторію польоту. На послідовності з 95 кадрів усі чотири конфігурації локалізували кожен кадр без втрати супроводу.
    \item Медіанна похибка локалізації на опорних кадрах становить $\approx 3$~\px{} для LoFTR і $\approx 6$~\px{} для LightGlue; не більше ніж 10~\px{} помиляється 94\,\% і 78--83\,\% кадрів відповідно. LoFTR точніший на більшості кадрів, і його перевага найпомітніша на складних ділянках (дах високої будівлі, дерева).
    \item LightGlue у 2,5--2,9 раза швидший за LoFTR у режимі трекінгу і допускає кешування ознак карти. Конфігурація LightGlue + MAGSAC++ --- єдина, що обробляє кадр не довше 1~с (0,78~с у середньому).
    \item MAGSAC++ скорочує час обробки на 17--28\,\% без статистично значущої зміни точності, тому його доцільно використовувати замість класичного RANSAC.
    \item Обидва методи стійкі до змін яскравості, контрасту, шуму та помірного розмиття, але не до повороту кадру понад $30$--$45^\circ$. Компенсація курсу за оцінкою кута з гомографії попередніх кадрів виявилася необхідною: без неї супровід втрачається на першому розвороті.
    \item Найефективнішим вдосконаленням є пошук у вікні карти за прогнозом: час кадру в режимі трекінгу скорочується в 2,9--3,4 раза для LightGlue (до 0,25--0,31~с) і в 5,1--6,7 раза для LoFTR (до 0,46--0,47~с) без статистично значущої зміни точності. Завдяки цьому точніший метод LoFTR також стає придатним для обробки в реальному часі на використаному обладнанні.
    \item Розбиття карти на 216 фрагментів $520 \times 520$ з перекриттям 33\,\% замість 384 фрагментів $512 \times 512$ з перекриттям 50\,\% пришвидшує глобальний пошук в 1,7--2,0 раза. Для LightGlue медіанна похибка при цьому зменшилася з 5,7--6,0 до 3,2--4,6~\px{}, однак виграш зосереджений на кількох складних кадрах і статистично незначущий; для LoFTR точність не змінилася.
    \item Штраф за відхилення масштабу, уточнення положення після глобального пошуку та згладжування траєкторії фільтром Калмана на досліджуваному польоті не підвищили точність. Уточнення для LightGlue навіть погіршило перший кадр через площину даху, а згладжування змінює медіанну похибку не більше ніж на 0,6~\px{}, що свідчить про систематичний, а не випадковий характер похибки.
    \item Вибір методу визначається вимогами застосування: LoFTR з вікном за прогнозом --- для вищої точності в реальному часі на GPU, LightGlue з вікном за прогнозом --- для найменшого часу обробки (0,25~с на кадр) і швидкого глобального пошуку завдяки кешуванню ознак карти.
    \item На другому польоті над полями (літній супутниковий знімок, листопадове відео) результати першого польоту не повторилися. Без компаса LightGlue визначив положення лише 7\,\% кадрів (медіанна похибка 5,2~м), а LoFTR на відрізку з орієнтирами --- жодного кадру. Хибних прив'язок при цьому не було. Над серединою полів зіставлення неможливе навіть за відомих положення й курсу, тому система локалізується лише за орієнтирами: дорогами, лісосмугами, трасою.
    \item Курс із компаса усуває невизначеність курсу під час глобального пошуку. LoFTR з компасом визначив положення 53--59\,\% кадрів відрізка з медіанною похибкою 2,2~м, і результат не погіршився за зміщення компаса на $10^\circ$ і шуму $5^\circ$. LightGlue з компасом визначив положення 43 кадрів замість 23 і знайшов нову ділянку польоту. Отже, компас --- доступний датчик, що суттєво підвищує надійність глобального пошуку, проте не замінює наявності орієнтирів на місцевості.
\end{enumerate}




\newpage
\section*{\centering ВИСНОВКИ}
\addcontentsline{toc}{section}{Висновки}

У магістерській роботі розроблено та досліджено систему абсолютної візуальної локалізації безпілотного літального апарата за ортофотопланом місцевості, яка визначає положення апарата без використання супутникової навігації на основі трансформерних методів зіставлення зображень LightGlue та LoFTR.

У теоретичній частині роботи проаналізовано підходи до візуальної навігації БпЛА в умовах відсутності або придушення сигналу GPS, принципи формування локальних ознак, згорткові кодувальники, механізми внутрішньої та перехресної уваги трансформерів, а також методи геометричної верифікації відповідностей. Розглянуто архітектури LightGlue, що зіставляє розріджені ключові точки з адаптивною глибиною обчислень, і LoFTR, що виконує бездетекторне щільне зіставлення за схемою <<від грубого до точного>>.

У практичній частині мовою Python розроблено програмну систему, яка забезпечує:
\begin{itemize}
    \item підготовку ортофотоплану: розбиття на фрагменти з перекриттям, перевірку цілісності, кешування ознак фрагментів;
    \item нормалізацію масштабу кадру і фрагмента карти та єдину систему координат у пікселях карти;
    \item зіставлення кадру з картою методами LightGlue і LoFTR через спільний інтерфейс, оцінювання гомографії методами RANSAC і MAGSAC++ та перевірку її правдоподібності;
    \item автоматичний пошук фрагмента карти для кадру без апріорної інформації про місце зйомки та обчислення координат апарата і контуру огляду камери;
    \item трекінг послідовності кадрів із прогнозом положення, кінематичною перевіркою, компенсацією курсу та відновленням після втрати супроводу;
    \item побудову незалежного еталона за ручною розміткою з автоматичним контролем її якості та обчислення метрик точності;
    \item вдосконалення: альтернативне розбиття карти, вікно карти за прогнозом, уточнення положення після глобального пошуку, штраф за відхилення масштабу, згладжування траєкторії фільтром Калмана та використання курсу з компаса;
    \item роботу з кількома наборами даних, географічну прив'язку карти та еталон за журналом супутникової навігації, який використовується лише для оцінювання;
    \item автоматизовані серії експериментів, що не змінюють основних результатів, і набір із 87 автоматичних тестів.
\end{itemize}

За результатами експериментального дослідження на першому польоті (95 кадрів над міською забудовою, зйомка взимку, ортофотоплан --- влітку) встановлено таке.
\begin{enumerate}
    \item Система автоматично знаходить відповідний фрагмент ортофотоплану та відстежує всю послідовність кадрів без утрати супроводу для всіх досліджених конфігурацій.
    \item Компенсація курсу за кутом, оціненим із гомографії попередніх кадрів, є необхідною умовою роботи системи: обидва методи втрачають працездатність при повороті кадру понад $30$--$45^\circ$, і без компенсації супровід утрачається на першому розвороті (74 з 95 кадрів).
    \item У основній конфігурації LoFTR точніший за LightGlue: медіанна похибка на 18 надійних опорних кадрах становить $\approx 3$~\px{} проти $\approx 6$~\px{}, а похибку не більше 10~\px{} мають 94\,\% і 78--83\,\% кадрів відповідно. LightGlue натомість у 2,5--2,9 раза швидший і допускає кешування ознак карти.
    \item Оцінювач MAGSAC++ скорочує час обробки на 17--28\,\% порівняно з RANSAC без статистично значущої зміни точності.
    \item Найефективнішим вдосконаленням є пошук у вікні карти за прогнозом: він скорочує час кадру в режимі трекінгу до 0,25--0,31~с для LightGlue і до 0,46--0,47~с для LoFTR (у 2,9--6,7 раза) без втрати точності. Отже, вимогу обробки кадру не довше ніж за 1~с на використаному ноутбуці з GPU GTX 1650 виконують обидва методи, зокрема точніший LoFTR.
    \item Розбиття карти на фрагменти $520 \times 520$ з перекриттям 33\,\% пришвидшує глобальний пошук в 1,7--2,0 раза; для LightGlue воно зменшило медіанну похибку до 3,2--4,6~\px{}, однак цей виграш статистично незначущий.
    \item Штраф за відхилення масштабу, уточнення положення після глобального пошуку та згладжування траєкторії точність на досліджуваному польоті не підвищили. Похибка має переважно систематичний характер: її визначають паралакс високих будівель, ділянки над деревами з сезонною зміною крон і невизначеність еталона.
    \item Результати курсової роботи щодо переваги LightGlue за кількістю відповідностей отримано за нерівних умов (різний масштаб вхідних зображень). За однакових умов LoFTR знаходить у 3,6--4 рази більше геометрично узгоджених відповідностей.
\end{enumerate}

Перевірка на другому польоті (335 кадрів над полями, листопадове відео, літній супутниковий знімок, еталон за даними GNSS) показала таке.
\begin{enumerate}[resume]
    \item Працездатність системи суттєво залежить від типу місцевості. Без компаса LightGlue визначив положення лише 7\,\% кадрів, а LoFTR на відрізку з орієнтирами --- жодного. Над однорідними полями з сезонною зміною покриву зіставлення неможливе навіть за відомих положення й курсу. Система не давала хибних прив'язок: усі видані положення мають похибку до 10~м (медіана 5,2~м).
    \item Курс із бортового компаса усуває невизначеність курсу під час глобального пошуку і є найдієвішим способом підвищити надійність на такій місцевості. З компасом LoFTR визначив положення 53--59\,\% кадрів відрізка з медіанною похибкою 2,2~м, навіть якщо компас мав зміщення $10^\circ$ і шум $5^\circ$; LightGlue визначив положення 43 кадрів замість 23.
\end{enumerate}

Таким чином, для систем візуальної локалізації БпЛА реального часу доцільно використовувати LoFTR із пошуком у вікні карти за прогнозом, якщо пріоритетом є точність, і LightGlue з тим самим вікном, якщо пріоритетом є найменший час обробки та швидкий глобальний пошук завдяки кешуванню ознак карти.

Для польотів над однорідною місцевістю систему доцільно доповнювати курсом із компаса, а ортофотоплан підбирати якомога ближчим до пори року польоту.

Основними напрямами подальшого розвитку є:
\begin{itemize}
    \item візуальна одометрія між сусідніми кадрами, щоб проходити ділянки без орієнтирів із виправленням накопиченого дрейфу на першому впізнаному орієнтирі;
    \item врахування непланарності сцени (висоти будівель) у моделі перетворення та в критеріях вибору гомографії;
    \item уточнення масштабу кадру під час польоту з урахуванням рельєфу;
    \item перевірка використання компаса з реальним магнітометром;
    \item перевірка системи на більшій кількості польотів, різних висотах, порах року й типах місцевості.
\end{itemize}

У результаті виконання магістерської роботи всі поставлені завдання виконано, а мету досягнуто.


\newpage
\section*{\centering СПИСОК ВИКОРИСТАНИХ ДЖЕРЕЛ}
\addcontentsline{toc}{section}{Список використаних джерел}

\begin{thebibliography}{99}

% Джерела кирилицею (за алфавітом)
\bibitem{kalman}
Kalman R. E. A New Approach to Linear Filtering and Prediction Problems // Journal of Basic Engineering. --- 1960. --- Vol. 82, No. 1. --- P. 35--45.

% Джерела латиницею (за алфавітом авторів)
\bibitem{magsac}
Barath D., Noskova J., Ivashechkin M., Matas J. MAGSAC++, a Fast, Reliable and Accurate Robust Estimator // 2020 IEEE/CVF Conference on Computer Vision and Pattern Recognition (CVPR). --- Seattle, WA, USA, 2020. --- P. 1301--1309.

\bibitem{opencv}
Bradski G. The OpenCV Library // Dr. Dobb's Journal of Software Tools. --- 2000. --- Vol. 25, No. 11. --- P. 120--123. --- URL: \url{https://opencv.org}

\bibitem{superpoint}
DeTone D., Malisiewicz T., Rabinovich A. SuperPoint: Self-Supervised Interest Point Detection and Description // 2018 IEEE/CVF Conference on Computer Vision and Pattern Recognition Workshops (CVPRW). --- Salt Lake City, UT, USA, 2018. --- P. 224--236.

\bibitem{esri}
Esri. World Imagery [Electronic resource] // ArcGIS Online. --- URL: \url{https://www.arcgis.com/home/item.html?id=10df2279f9684e4a9f6a7f08febac2a9}

\bibitem{ransac}
Fischler M. A., Bolles R. C. Random sample consensus: a paradigm for model fitting with applications to image analysis and automated cartography // Communications of the ACM. --- 1981. --- Vol. 24, No. 6. --- P. 381--395.

\bibitem{hartley}
Hartley R., Zisserman A. Multiple View Geometry in Computer Vision. --- 2nd ed. --- Cambridge: Cambridge University Press, 2004. --- 655 p.

\bibitem{resnet}
He K., Zhang X., Ren S., Sun J. Deep Residual Learning for Image Recognition // 2016 IEEE Conference on Computer Vision and Pattern Recognition (CVPR). --- Las Vegas, NV, USA, 2016. --- P. 770--778.

\bibitem{cnn_fpn}
Lin T.-Y., Dollár P., Girshick R., He K., Hariharan B., Belongie S. Feature Pyramid Networks for Object Detection // 2017 IEEE Conference on Computer Vision and Pattern Recognition (CVPR). --- Honolulu, HI, USA, 2017. --- P. 2117--2125.

\bibitem{lightglue}
Lindenberger P., Sarlin P.-E., Pollefeys M. LightGlue: Local Feature Matching at Light Speed // 2023 IEEE/CVF International Conference on Computer Vision (ICCV). --- Paris, France, 2023. --- P. 17627--17638.

\bibitem{sift}
Lowe D. G. Distinctive Image Features from Scale-Invariant Keypoints // International Journal of Computer Vision. --- 2004. --- Vol. 60, No. 2. --- P. 91--110.

\bibitem{classic_methods}
Mikolajczyk K., Schmid C. A Performance Evaluation of Local Descriptors // IEEE Transactions on Pattern Analysis and Machine Intelligence. --- 2005. --- Vol. 27, No. 10. --- P. 1615--1630.

\bibitem{pytorch}
Paszke A., Gross S., Massa F. et al. PyTorch: An Imperative Style, High-Performance Deep Learning Library // Advances in Neural Information Processing Systems (NeurIPS). --- 2019. --- Vol. 32. --- P. 8024--8035.

\bibitem{qgis}
QGIS Development Team. QGIS Geographic Information System [Electronic resource]. --- Open Source Geospatial Foundation. --- URL: \url{https://qgis.org}

\bibitem{rts}
Rauch H. E., Tung F., Striebel C. T. Maximum Likelihood Estimates of Linear Dynamic Systems // AIAA Journal. --- 1965. --- Vol. 3, No. 8. --- P. 1445--1450.

\bibitem{kornia}
Riba E., Mishkin D., Ponsa D., Rublee E., Bradski G. Kornia: an Open Source Differentiable Computer Vision Library for PyTorch // 2020 IEEE Winter Conference on Applications of Computer Vision (WACV). --- Snowmass Village, CO, USA, 2020. --- P. 3663--3672.

\bibitem{orb}
Rublee E., Rabaud V., Konolige K., Bradski G. ORB: An efficient alternative to SIFT or SURF // 2011 International Conference on Computer Vision (ICCV). --- Barcelona, Spain, 2011. --- P. 2564--2571.

\bibitem{visual_nav}
Scaramuzza D., Fraundorfer F. Visual Odometry [Tutorial] // IEEE Robotics \& Automation Magazine. --- 2011. --- Vol. 18, No. 4. --- P. 80--92.

\bibitem{loftr}
Sun J., Shen Z., Wang Y., Bao H., Zhou X. LoFTR: Detector-Free Local Feature Matching with Transformers // 2021 IEEE/CVF Conference on Computer Vision and Pattern Recognition (CVPR). --- Nashville, TN, USA, 2021. --- P. 8918--8927.

\bibitem{attention_all_need}
Vaswani A., Shazeer N., Parmar N., Uszkoreit J., Jones L., Gomez A. N., Kaiser L., Polosukhin I. Attention Is All You Need // Advances in Neural Information Processing Systems (NeurIPS). --- 2017. --- Vol. 30. --- P. 5998--6008.

\bibitem{exg}
Woebbecke D. M., Meyer G. E., Von Bargen K., Mortensen D. A. Color Indices for Weed Identification Under Various Soil, Residue, and Lighting Conditions // Transactions of the ASAE. --- 1995. --- Vol. 38, No. 1. --- P. 259--269.

\end{thebibliography}

\end{document}